\chapter[Langages rationnels]{Automates finis et langages rationnels}
\label{chp.auto}

\minitoc

\lettrine{P}{our} aborder la calculabilité, plusieurs approches existent. L'une
d'entre elles est de considérer directement les formalismes de fonctions
calculables. Cette approche est raisonnable puisqu'elle permet d'aller à
l'essentiel, mais cela donne une vision assez restrictive des modèles de calculs
plus faibles existants, et qui sont souvent très utiles.

C'est pourquoi nous faisons le choix dans cet ouvrage (où nous ne sommes limités
ni par le temps d'un cours ni par l'espace d'un livre imprimé) de traiter, avant
les formalismes de fonctions calculables, les formalismes liés à la hiérarchie
de Chomsky. Nous pensons en particulier que les automates permettent d'illuminer
profondément le fonctionnement intuitif d'une machine de Turing, plus complexe
mais basé sur des notions très proches.

Ce chapitre couvre donc les bases de l'étude des langages rationnelss et des
automates finis. Cette classe de langages peut être décrite par de nombreux
formalismes, mais nous nous contenterons des automates finis, des expressions
rationnelless et des monoïdes.

Ce chapitre commence par la base de la théorie des langages~: les définitions
de ce qu'est un alphabet, un mot, un langage\ldots Nous abordons ensuite la
classe des langages rationnels en considérant les formalismes que sont les
automates et les expressions rationnelles, dont on prouve leur équivalence.
Nous étudions après cela les théorèmes de structure sur ces
langages~: le lemme de l'étoile et le théorème de Myhill-Nérode, ainsi que ses
conséquences. Nous traitons en dernier le cas de l'équivalence des langages
rationnels avec les langages reconnus par des monoïdes, pour ainsi revenir sur
un aspect plus algébrique, proche du début du chapitre.

\section{Alphabet, mot, langage}

Pour étudier de façon plus systématique les langages, nous étudions d'abord de
l'algèbre avec la structure de monoïde, qui est intimement liée à celle
d'ensemble de mots. Tout d'abord, nous donnons une description concrète de ce
qu'est un mot et un langage, avant d'aborder l'aspect théoriques de l'étude des
monoïdes.

\subsection{Premières définitions}

Définissons donc la notion d'alphabet. Intuitivement, l'alphabet désigne
l'ensemble des symboles que nous utilisons pour écrire. Pour avoir un traitement
uniforme, on considère que des symboles comme \og ( \fg font aussi partie de
l'alphabet. C'est alors un objet particulièrement quelconque~: on peut imaginer
écrire avec à peu près n'importe quel ensemble de symboles. D'où la définition
semblant inutile d'un alphabet~: parler d'alphabet permet surtout de situer le
contexte d'étude.

\begin{definition}[Alphabet]
  Un alphabet $\ABA$ est un ensemble fini non vide.
\end{definition}

L'intérêt d'un alphabet est avant tout de pouvoir écrire des mots avec. La
notion de mot est la notion finie par excellence~: un mot est donc simplement
une suite finie de lettres.

\begin{definition}[Mot]
  Un mot $\varwd$ sur un alphabet $\ABA$ est une suite finie à valeurs dans
  $\ABA$, c'est-à-dire $\varwd = (\varwd_0,\ldots,\varwd_{n-1})\in \ABA^n$ pour
  un certain entier $n\in\bN$.

  On note $\toStA$ l'ensemble des mots sur $\ABA$~:
  \[\toStA \defeq \bigcup_{n\in \bN} \ABA^n\]
\end{definition}

\begin{notation}
  Pour tout mot $\varwd\in\ABA^n$, on note $\varwd_i$ pour la
  \ordinalnumeralfeminin{$i$} lettre, c'est-à-dire l'image de $\varwd$ par la
  projection $\projP{i}{} : \ABA^n \to \ABA$.

  Pour tout mot $\varwd\in\toStA$, on note $\wul\in\bN$ l'entier $n$ tel que
  $\varwd\in \ABA^n$. On définit aussi, pour $\varwd\in\toStA$ et
  $\varla\in\ABA$, l'entier $\wul_\varla \in \bN$ comme le nombre d'occurrences
  de $\varla$ dans $\wul$.
\end{notation}

\begin{remark}
  On aurait aussi pu définir $\toStA$ comme l'ensemble inductif (au sens du
  \cref{chp.induction}) des listes d'éléments de $\ABA$. Les deux formalismes
  sont équivalents, mais la théorie des automates et plus généralement celle des
  langages formels utilisent habituellement plutôt le formalisme des suites
  finies. On peut aussi remarquer une inversion dans le sens de lecture. Dans
  le formalisme des listes, l'élément en tête d'une liste $a\Cons \varlist$ est
  à gauche, tandis que la lecture d'un mot $\varwd \in \toStA$ se fait de gauche
  à droite.
\end{remark}

Remarquons qu'une suite vide est une suite finie~: on peut donc définir le mot
vide, qu'on notera $\Ewd$.

On peut désormais définir un langage~: c'est un ensemble de mots.

\begin{definition}[Langage]
  Soit $\ABA$ un alphabet. Un langage $\varLL$ sur $\ABA$ est une partie
  $\varLL\subseteq\toStA$.
\end{definition}

Parlons dès maintenant de stabilité, qui est une notion importante dans le
contexte de l'algèbre~: elle permet de décrire une structure dans laquelle
une opération donnée est interne à la structure.

\begin{definition}[Stabilité]
  Soit $\ABA$ un alphabet et $\varOpe$ une opération binaire sur $\toStA$. On
  dit qu'un langage $\varLL$ sur $\ABA$ est stable par $\varOpe$ si
  \[\forall \varwd,\varwv\in \varLL, \varwd\varOpe \varwv \in \varLL\]
\end{definition}

Une autre façon de parler de stabilité est de commencer par étendre une
opération binaire sur les mots en une opération sur les langages.

\begin{definition}[Extension sur les langages]
  Soit $\ABA$ un alphabet et $\varOpe$ une opération binaire sur $\toStA$. On
  étend l'opération $\varOpe$ sur les langages sur $\ABA$ par, pour tous
  langages $\varLL,\varLLB$~:
  \[\varLL\varOpe \varLLB \defeq
  \compr{\varwd\varOpe\varwv}{\varwd \in \varLL, \varwv \in \varLLB}\]
\end{definition}

\begin{exercise}
  Montrer qu'un langage $\varLL$ est stable par une opération $\varOpe$ si et
  seulement si $\varLL \varOpe \varLL \subseteq \varLL$.
\end{exercise}

Cette notion de stabilité se généralise très bien à des classes de langages,
dont il convient donc d'introduire la définition.

\begin{definition}[Classe de langages]
  Soit un alphabet $\ABA$. On appelle classe de langages une partie
  $\varCC\subseteq \powerset(\toStA)$ (ses éléments sont donc des langages).
\end{definition}

\begin{definition}[Stabilité d'une classe de langages]
  Soit $\ABA$ un alphabet et $\varOpe$ une opération binaire sur
  $\powerset(\ABA)$. On dit qu'une classe de langages $\varCC$ sur $\ABA$ est
  stable par $\varOpe$ si
  \[\forall \varLL,\varLLB\in\varCC, \varLL\varOpe \varLLB\in\varCC\]
\end{definition}

\begin{remark}
  Si on dispose d'une opération binaire $\varOpe$ sur les langages, il est
  possible de l'étendre sur les classes en faisant une construction analogue
  à celle étendant les opérations sur les mots en des opérations sur les
  langages. On peut ainsi passer d'une opération sur les mots en une opération
  sur les classes de langages, en effectuant les deux extensions à la suite.
\end{remark}

Pour le reste de cette sous-section, nous allons voir des exemples de langages,
d'opérations et de classes de langages.

\begin{example}
  Nous allons dès maintenant définir l'alphabet qui nous sera le plus utile~:
  \[\btwo \defeq \intB\]

  Voici quelques langages sur $\btwo$ ou un alphabet $\ABA$ quelconque~:
  \begin{itemize}
  \item le langage des mots commençant par $0$~:
    \[\varLL_0 \defeq \compr{\varwd\in\toStA}{\varwd_0 = 0}\]
  \item le langage des mots ayant une taille paire~:
    \[\varLL_{\pair} \defeq \compr{\varwd\in\toStA}{\wul\in 2\bN}\]
  \item le langage des mots dont le début et la fin sont la même lettre~:
    \[\varLL_{\mathrm{debut}=\mathrm{fin}}\defeq
    \compr{\varwd\in\toStA}{\varwd_0 = \varwd_{\wul - 1}}\]
  \item pour un langage $\varLL$, le langage $\mirL{\varLL}$ des mots pris dans
    l'autre sens~:
    \[\mirL\varLL\defeq
    \compr{\varwd\in\toStA}{\varwd_{\wul-1}\ldots\varwd_0\in\varLL}\]
  \item étant donnés deux langages $\varLL$ et $\varLLB$, l'intersection des
    deux langages $\varLL \cap \varLLB$.
  \item étant donnés deux langages $\varLL$ et $\varLLB$, l'union des deux
    langages $\varLL\cup\varLLB$.
  \end{itemize}
\end{example}

Nous allons définir une nouvelle opération, fortement liée à la structure de
$\toStA$, qui nous permettra de donner plus d'exemples de langages.

\begin{definition}[Concaténation]
  Soit un alphabet $\ABA$. On définit l'opération de concaténation
  $\makeFunName{\tok\wdPt\tok}{\toStA\times\toStA}{\toStA}$ comme le fait,
  étant donnés deux mots $\varwd,\varwv\in\toStA$, de lire les deux mots l'un à
  la suite de l'autre. Formellement, on définit
  \[\varwd\wdPt\varwv \defeq
  (\varwd_0,\ldots,\varwd_{\wul-1},\varwv_0,\ldots,\varwd_{\wvl-1})\]

  Pour tout $n\in \bN$, on notera $\varwd^n$ pour l'itération de l'opération
  $\wdPt$~:
  \[\begin{cases}
  \varwd^0 \defeq \Ewd \\
  \forall n \in \bN, \varwd^{n+1} = \varwd^n \wdPt \varwd
  \end{cases}\]
\end{definition}

\begin{remark}
  Le mot vide, $\Ewd$, est neutre pour $\wdPt$~: $\varwd\wdPt\Ewd = \varwd$
  et $\Ewd \wdPt \varwd = \varwd$.
\end{remark}

\begin{notation}
  Il nous arrivera régulièrement de ne pas écrire l'opération de concaténation.
  La notation $\varwd\varwv$ où $\varwd,\varwv\in\toStA$, doit donc se lire
  $\varwd\wdPt\varwv$.
\end{notation}

\begin{example}
  Voici plusieurs autres langages possibles sur un alphabet $\ABA$~:
  \begin{itemize}
  \item pour deux langages $\varLL,\varLLB$, le langage concaténé des deux
    langages~:
    \[\varLL\wdPt\varLLB\defeq
    \compr{\varwd\wdPt\varwv}{\varwd\in\varLL,\varwv\in \varLLB}\]
  \item pour un langage $\varLL$, l'itération de concaténation sur ce langage~:
    \[\begin{cases}
    \varLL^{\wdPt 0} \defeq \{\Ewd\}\\
    \forall n \in \bN, \varLL^{\wdPt(n+1)} \defeq \varLL^{\wdPt n}\wdPt\varLL
    \end{cases}\]
  \item pour un langage $\varLL$, le langage des mots dont toutes les
    concaténations sont encore dans $\varLL$~:
    \[\varLL_{\stab\mathord{\wdPt}}\defeq
    \compr{\varwd\in\toStA}{\forall n \in \bN, \varwd^n \in \varLL}\]
  \item pour un langage $\varLL$, le langage des itérations de mots de
    $\varLL$~:
    \[\toSt{\varLL} \defeq
    \compr{\varwd_0\wdPt\cdots\wdPt\varwd_{n-1}}{n \in \bN,
      \forall i\in \intervInt{n}, \varwd_i \in \varLL}\]
  \end{itemize}
\end{example}

\begin{exercise}
  Montrer que pour tout langage $\varLL$ sur un alphabet $\ABA$, on a l'égalité
  suivante~:
  \[\toSt{\varLL} = \bigcup_{n \in \bN} \varLL^{\wdPt n}\]
\end{exercise}

\begin{exercise}
  Montrer que pour tout langage $\varLL$ sur un alphabet $\ABA$, on a
  $\Ewd \in \varLL^{\wdPt}$ et $\varLL^{\wdPt}$ est stable par $\wdPt$.
\end{exercise}

\begin{example}
  Donnons deux exemples de classes de langages sur un alphabet $\ABA$~:
  \begin{itemize}
  \item la classe des langages contenant $\Ewd$~:
    \[\varCC_\Ewd \defeq \compr{\varLL \subseteq \toStA}{\Ewd\in\varLL}\]
  \item la classe des langages stables par $\wdPt$~:
    \[\varCC_{\wdPt} \defeq
    \compr{\varLL \subseteq\toStA}
          {\forall \varwd,\varwv\in\varLL, \varwd\wdPt\varwv \in \varLL}
    \]
  \end{itemize}
\end{example}

\subsection{Aspect algébrique des langages}

Nous nous intéressons maintenant à l'aspect algébrique des langages. L'objectif
de cette sous-section est d'aboutir à la caractérisation de l'ensemble des mots
comme construction universelle (celle du monoïde libre). La
\cref{def.monoid.ci} permet déjà de deviner la définition d'un monoïde~: on
demande simplement moins de conditions.

\begin{definition}[Monoïde]
  Un monoïde est un ensemble $\varMMo$ muni d'une opération binaire
  $\tok\opM\tok : \varMMo \times \varMMo \to \varMMo$ et d'un élément neutre
  $\neuMM$, c'est-à-dire un élément $\neuMM\in \varMMo$ tel que
  \[\forall x\in\varMMo, x\opM\neuMM = \neuMM\opM x = x\]
  et tel que $\opM$ est associative, c'est-à-dire
  \[\forall x,y,z\in\varMMo, x\opM (y\opM z) = (x\opM y) \opM z\]
\end{definition}

\begin{example}
  L'exemple le plus évident de monoïde est $(\bN,+,0)$.
\end{example}

\noindent
Définissons la notion de morphisme associée à la structure de monoïde.

\begin{definition}[Morphisme de monoïde]
  Soient $(\varMMo,\opM,\neuMM)$ et $(\varMNo,\opN,\neuMN)$ deux monoïdes. Un
  morphisme entre ces deux monoïdes est une fonction $f : \varMMo \to \varMNo$
  vérifiant les deux conditions
  \begin{equation*}
    \forall x,y\in \varMMo, f(x\opM y) = f(x)\opN f(y)
  \end{equation*}
  \begin{equation*}
    f(\neuMM) = \neuMN
  \end{equation*}
\end{definition}

Remarquons que contrairement au cas des groupes, les morphismes de monoïdes ont
besoin de stabiliser l'élément neutre.

\begin{exercise}
  Soit $\ABA$ un alphabet, vérifier que $(\toStA, \wdPt,\Ewd)$ est un monoïde.
\end{exercise}

Le cas de $\toStA$ généralise le cas de $\bN$~: il est en effet possible de
considérer un entier dans sa représentation unaire (comme une
succession d'applications de la fonction $n \mapsto n + 1$ à $0$) comme un mot
sur l'alphabet $\{*\}$ où $*$ est un objet quelconque. La seule information
que contient un tel mot est alors sa longueur.

L'étude des monoïdes permet donc d'inclure l'étude de la structure de
$\toStA$. Nous allons maintenant voir qu'en fait, $\toStA$ peut être construit
comme la structure de monoïde la plus simple engendrée par $\ABA$. Cela se
traduit par la propriété universelle du monoïde libre.

\begin{theorem}[Propriété universelle du monoïde libre]\label{thm.PU.monoid}
  Soit $\ABA$ un alphabet, $(\varMMo,\opM,\neuMM)$ un monoïde et
  $f :\ABA\to\varMMo$ une fonction (quelconque). Alors il existe un unique
  morphisme de monoïde $\relev{f}$ entre $\toStA$ et $\varMMo$ tel que
  $\relev{f}(x) = f(x)$ pour tout $x \in \ABA$.
\end{theorem}

\begin{proof}
  On procède par analyse-synthèse~:
  \begin{itemize}
  \item supposons que $\relev{f}$ existe et vérifie les conditions décrites.
    Alors pour tout mot $\varwd = \varwd_0\ldots \varwd_{n-1}$, on peut vérifier
    par récurrence sur $n$ que
    $\relev{f}(u) = f(\varwd_0)\opM \cdots \opM f(\varwd_{n-1})$, où le cas où
    $\varwd$ est de longueur nulle correspond à $\relev{f}(\Ewd) = \neuMM$. La
    fonction ainsi définie est unique, étant donnée la construction.
  \item en considérant la fonction proposée précédemment, on voit d'abord que
    $\relev{f}(x) = f(x)$ pour $x\in \ABA$, et que si $\varwd=\varwv\wdPt\varww$
    alors $\relev{f}(\varwd) = \relev{f}(\varwv)\opM \relev{f}(\varww)$, où
    l'associativité est nécessaire pour pouvoir couper le produit au milieu sans
    changer le résultat.
  \end{itemize}
  On voit donc qu'un seul candidat est possible pour la fonction donnée, et que
  ce candidat vérifie la proposition attendue.
\end{proof}

Remarquons qu'en choisissant une fonction $f : \ABA \to \varMMo$ sans propriété
particulière, on peut potentiellement atteindre avec $\relev{f}$ des valeurs qui
ne sont pas atteintes par $f$ directement. Par exemple si on prend la fonction
\[\makeFun{f}{\ABA}{\bN}{\varla}{1}\]
alors l'image d'un mot de taille $n$ est $n$. De cette façon, on vient de voir
que l'ensemble des mots de taille $n$ peut être décrit à partir du monoïde
$(\bN,+,1)$ en considérant une partie du monoïde et une fonction
$f : \ABA \to \bN$.

Plus généralement, on définit ce qu'est le langage induit par une fonction
depuis un alphabet vers un monoïde possédant une partie distinguée.

\begin{definition}
  Soit $\ABA$ un alphabet, $(\varMMo,\opM,\neuMM)$ un monoïde,
  $h : \ABA \to \varMMo$ une fonction et $S\subseteq\varMMo$. On dit que
  $\varLL(\ABA,\varMMo,\opM,\neuMM,h,S)$ est le langage induit par
  $(\ABA,\varMMo,\opM,\neuMM,h,S)$, qui est défini par~:
  \[\varLL(\ABA,\varMMo,\opM,\neuMM,h,S)\defeq
  \compr{\varwd\in\toStA}{\relev{h}(\varwd)\in S}\]
\end{definition}

\noindent
En reprenant l'exemple de $\varla \mapsto 1$, on voit qu'en prenant $S = \{n\}$,
on induit le langage $\ABA^n$.

\section{Automate fini et langage reconnaissable}

\noindent
Tournons-nous maintenant vers les automates finis et les langages qu'ils
reconnaissent.

\subsection{Définition d'un automate}

Pour présenter un automate, le meilleur moyen est de commencer par une approche
graphique~: un automate se visualise comme une machine effectuant des opérations
élémentaires sur des mots. Dans notre premier cas, l'opération est
particulièrement simple puisque la machine se contente de lire le mot. Un
automate fini est une machine lisant linéairement un mot (elle lit chaque lettre
l'une après l'autre, sans pouvoir faire demi-tour). A chaque lettre lue,
l'automate change (ou non) d'état, parmi un ensemble fini d'états~: le mot lu
est accepté si la lecture aboutit à un certain état considéré comme acceptant.

Le formalisme derrière cette définion est lourd, mais nécessaire, nous donnons
donc la définition formelle d'un automate.

\begin{definition}[Automate fini déterministe]
  Soit un alphabet $\ABA$. Un automate fini déterministe $\varAuA$ sur $\ABA$
  est un quadruplet $(\varQQ,\varqZ,\trA,\varAF)$ où~:
  \begin{itemize}
  \item $\varQQ$ est un ensemble fini appelé ensemble des états de $\varAuA$,
  \item $\varqZ \in \varQQ$ est appelé l'état initial,
  \item $\makeFunNamePartial{\trA}{\varQQ\times \ABA}{\varQQ}$ est une fonction
    partielle appelée la fonction (ou table) de transitions,
  \item $\varAF \subseteq \varQQ$ est appelé l'ensemble des états terminaux.
  \end{itemize}

  On notera $\auto(\ABA)$ l'ensemble des automates sur $\ABA$.
\end{definition}

Le processus de lecture décrit plus haut, en reprenant les notations
introduites, consiste à prendre un mot $\varwd = \varwd_0\ldots \varwd_n$ et
calculer $\trA(\trA(\cdots\trA(\varqZ,\varwd_0)\cdots),\varwd_n)$ puis vérifier
que cette valeur appartient à $\varAF$.

\begin{definition}[Fonction de transition étendue]
  Soit un automate fini déterministe $\varAuA$ sur un alphabet $\ABA$ et $\trA$
  sa fonction de transition. On appelle fonction de transition étendue la
  fonction $\makeFunNamePartial{\trAS}{\varQQ\times\toStA}{\varQQ}$ définie par
  induction sur $\toStA$ par~:
  \[\begin{cases}
  \trAS(\varq,\Ewd) = \varq \\
  \trAS(\varq,\varwd\wdPt \varla) = \trA(\trAS(\varq,\varwd),\varla)
  \end{cases}\]
  pour tout état $q \in \varQQ$, mot $\varwd \in \toStA$ et lettre
  $\varla \in \ABA$.
\end{definition}

\begin{definition}[Langage accepté par un automate]
  Soit $\varAuA$ un automate sur un alphabet $\ABA$. On dit que $\varAuA$
  accepte (ou reconnait) le mot $\varwd \in \toStA$ si
  $\trAS(\varqZ,\varwd) \in \varAF$, ce qu'on note $\varAuA \satisf \varwd$. On
  définit le langage reconnu par $\varAuA$ par~:
  \[\varLL_\varAuA \defeq \compr{\varwd \in \toStA}{\varAuA\satisf\varwd}\]
\end{definition}

Maintenant que les définitions ont été introduites, donnons une intuition
graphique de ce qu'il se passe (intuition largement indispensable pour
comprendre les automates). Pour représenter un automate $\varAuA$, la convention
est de dessiner un graphe donc les sommets sont les états (l'ensemble $\varQQ$),
et où $\trA$ est représentée par des arcs dirigés et étiquetés par $\ABA$. Par
exemple, entre deux états $\varq$ et $\varq'$, s'il existe $\varla$ tel que
$\trA(\varq,\varla) = \varq'$, alors on aura une flèche
$\varq \xrightarrow{\varla} \varq'$. L'état initial est représenté avec une
flèche entrante, et les états terminaux sont représentés par des doubles
cercles.

Prenons par exemple un premier automate $\varAuA_0$ sur $\btwo$ donné dans la
\cref{fig.A0}.

\begin{figure}[t]
  \centering
  \begin{tikzpicture}[node distance = 2cm, on grid, auto]
    \node[state, initial] (q_0) {$\varqP{0}$};
    \node[state,accepting] (q_1) [right = of q_0] {$\varqP{1}$};
    \node[state] (q_2) [above right = of q_1] {$\varqP{2}$};
    \node[state, accepting] (q_3) [below right = of q_1] {$\varqP{3}$};
    \path[->,>=latex]
    (q_0) edge node {0,1} (q_1)
    (q_1) edge node {1} (q_2)
    edge node {0} (q_3)
    (q_2) edge node {0} (q_3)
    edge [loop above] node {$1$} ()
    (q_3) edge node {1} (q_0);
  \end{tikzpicture}
  
  \vspace{0.2cm}
  \hrule

  \vspace{0.2cm}
  \caption{Automate $\varAuA_0$ sur l'alphabet $\btwo$}
  \label{fig.A0}
  
\end{figure}

\begin{figure}[t]
  \centering
  \hrule

  \vspace{0.2cm}
  \begin{tikzpicture}[node distance = 2cm, on grid, auto,
      scale=0.65, transform shape]
    \node[state, initial, color = red] (q_0) {$\varqP{0}$};
    \node[state,accepting] (q_1) [right = of q_0] {$\varqP{1}$};
    \node[state] (q_2) [above right = of q_1] {$\varqP{2}$};
    \node[state, accepting] (q_3) [below right = of q_1] {$\varqP{3}$};
    \path[->,>=latex]
    (q_0) edge node {0,1} (q_1)
    (q_1) edge node {1} (q_2)
    edge node {0} (q_3)
    (q_2) edge node {0} (q_3)
    edge [loop above] node {$1$} ()
    (q_3) edge node {1} (q_0);

    \node (a) [right = of q_1] {};
    \node (b) [right = of a] {};
    \path[|->, >=latex] (a) edge node {$\trA(\varqZ,0)$} (b);

    \node[state, initial] (p_0) [right = of b] {$\varqP{0}$};
    \node[state,accepting, color = red] (p_1) [right = of p_0] {$\varqP{1}$};
    \node[state] (p_2) [above right = of p_1] {$\varqP{2}$};
    \node[state, accepting] (p_3) [below right = of p_1] {$\varqP{3}$};
    \path[->,>=latex]
    (p_0) edge node {0,1} (p_1)
    (p_1) edge node {1} (p_2)
    edge node {0} (p_3)
    (p_2) edge node {0} (p_3)
    edge [loop above] node {$1$} ()
    (p_3) edge node {1} (p_0);

    \node (aa) [right = of p_1] {};
    \node (bb) [right = of aa] {};
    \path[|->, >=latex] (aa) edge node {$\trA(\varqP{1},0)$} (bb);

    \node[state, initial] (r_0) [ right = of bb] {$\varqZ$};
    \node[state,accepting] (r_1) [right = of r_0] {$\varqP{1}$};
    \node[state] (r_2) [above right = of r_1] {$\varqP{2}$};
    \node[state, accepting, color = red] (r_3) [below right = of r_1]
         {$\varqP{3}$};
    \path[->,>=latex]
    (r_0) edge node {0,1} (r_1)
    (r_1) edge node {1} (r_2)
    edge node {0} (r_3)
    (r_2) edge node {0} (r_3)
    edge [loop above] node {$1$} ()
    (r_3) edge node {1} (r_0);
  \end{tikzpicture}
  
  \vspace{0.2cm}
  \hrule

  \vspace{0.2cm}
  \caption{Lecture du mot $00$ par l'automate $\varAuA_0$}
  \label{fig.ex.reco}
\end{figure}

En lisant simplement ce dessin, on peut voir que l'automate reconnaît, par
exemple, $00$ (la lecture de ce mot est donnée en \cref{fig.ex.reco}). Cette
représentation a l'avantage, en plus de la lisibilité, de donner l'ensemble des
informations nécessaires pour retrouver l'objet mathématique représenté associé.

On a ainsi construit une classe de langages~: la classe des langages
reconnaissables.

\begin{definition}[Langage reconnaissable]
  Soit $\ABA$ un alphabet. On appelle classe des langages reconnaissbles sur
  $\ABA$ l'ensemble des langages reconnus par les automates sur $\ABA$~:
  \[\Reco(\ABA) \defeq \compr{\varLL_{\varAuA}}{\varAuA\in\auto(\ABA)}\]
\end{definition}

\begin{exercise}
  Trouver le langage reconnu par l'automate $\varAuA_1$ sur $\btwo$
  donné en \cref{fig.auto.exo1} (et prouver que le langage reconnu est bien
  celui attendu)~:
  \begin{figure}[t]
    \centering
    \begin{tikzpicture}[node distance = 2cm, on grid, auto]
      \node[state, initial,accepting] (q_0) {$\varqP{0}$};
      \node[state] (q_1) [right = of q_0] {$\varqP{1}$};
      \path[->,>=latex]
      (q_0) edge [bend left] node {0} (q_1)
      edge [loop above] node {1} (q_0)
      (q_1) edge [bend left] node {0} (q_0)
      edge [loop above] node {1} (q_1);
    \end{tikzpicture}
  
  \vspace{0.2cm}
  \hrule

  \vspace{0.2cm}
  \caption{Automate $\varAuA_1$}
  \label{fig.auto.exo1}
  \end{figure}
\end{exercise}

Cette première classe de langages d'intérêt est en fait le centre de l'étude de
ce chapitre. On le verra, mais plusieurs définitions de cette classe existent,
par des formalismes très différents, ce qui rend cette classe de langages assez
remarquable.

L'un de ces formalismes a déjà été présenté~: on verra que les langages associés
à des monoïdes sont en fait des langages reconnus par des automates.

\subsection{Langage rationnel}

Avant de donner notre dernier formalisme, décrivons une autre classe de
langages, qui s'avérera en fait être la même classe.

\begin{definition}[Langage rationnel]
  La classe des langages rationnels sur un alphabet $\ABA$ est la plus petite
  classe $\varCR$ de langages sur $\ABA$ telle que~:
  \begin{itemize}
  \item $\vide \in \varCR$
  \item pour tout $\varla \in \ABA$, $\{\varla\}\in \varCR$
  \item si $\varLL,\varLLB \in \varCR$ alors $\varLL\cup \varLLB \in \varCR$
  \item si $\varLL,\varLLB \in \varCR$ alors $\varLL\wdPt \varLLB \in \varCR$
  \item si $\varLL \in \varCR$ alors $\toSt{\varLL} \in \varCR$.
  \end{itemize}

  On note $\Ratio(\ABA)$ cette classe $\varCR$.
\end{definition}

\begin{remark}
  Le langage $\{\Ewd\}$ est rationnel, car il peut s'écrire $\toSt{\vide}$. On
  peut donc ajouter si on le souhaite le mot vide tout en gardant la
  rationnalité d'un langage.
\end{remark}

Cette définition est une définition par induction, comme on en a rencontré
de nombreuses fois depuis le début de ce livre. Cependant, il n'est pas
particulièrement facile de manipuler un tel ensemble de langages, en particulier
pour représenter un langage de façon efficace~: devoir décrire le procédé par
lequel on obtient le langage est, pour l'instant, fastidieux.

Heureusement, il existe un formalisme permettant exactement de décrire comment
on obtient un langage rationnel à partir de ces opérations inductives, c'est ce
qu'on appelle les expressions rationnelles.

\begin{definition}[Expression rationnelle]
  Soit un alphabet $\ABA$. On définit l'ensemble des expressions rationnelles
  par la grammaire suivante~:
  \[e,e' \inddefeq \vide\mid \varla \mid e + e' \mid e e' \mid \toSt{e}\]
  où $\varla \in \ABA$. On note $\Regex(\ABA)$ l'ensemble des expressions
  rationnelles sur $\ABA$.
\end{definition}

\begin{remark}
  On appelle aussi ces expressions et langages des expressions régulières et des
  langages réguliers, d'où le terme \foreignexpr{regular expression} donnant
  lieu à \foreignexpr{regex}.
\end{remark}

La correspondance est assez naturelle~: une expression de la forme $\varla$
représente le langage $\{\varla\}$, $\vide$ le langage vide et les trois
opérations représentent les opérations par lesquelles est stable la classe des
langages rationnels.

On définit donc une fonction d'interprétation des expressions rationnelles. De
plus, on ajoute $\Ewd$ dans nos expressoins rationnelles, par souci de
lisbilité.

\begin{definition}[Interprétation d'une expression rationnelle]
  On définit la fonction $\makeFunName{\Val}{\Regex(\ABA)}{\Ratio(\ABA)}$ par
  induction~:
  \begin{itemize}
  \item $\Val(\vide) = \vide$
  \item $\Val(\varla) = \{\varla\}$
  \item $\Val(e+e') = \Val(e) \cup \Val(e')$
  \item $\Val(ee') = \Val(e) \wdPt \Val(e')$
  \item $\Val(\toSt{e}) = \toSt{(\Val(e))}$
  \item $\Val(\Ewd) = \{\Ewd\}$
  \end{itemize}
\end{definition}

\begin{exercise}
  Montrer que $\Val$ est une surjection.
\end{exercise}

\begin{exercise}
  Trouver le langage interprété par l'expression rationnelle suivante sur
  $\btwo$~:
  \[\toSt{(0\toSt{1} 0)}\]
\end{exercise}

\begin{remark}
  Cette surjection n'est pas pour autant injective, et ce même en retirant
  $\Ewd$ des expressions rationnelles. Par exemple, $\varla + \varla$ a la même
  interprétation que $\varla$.
\end{remark}

On a ainsi une façon plus simple de manipuler la classe des langages rationnels.
Notre première équivalence va donc consister en deux étapes~:
\begin{itemize}
\item montrer que les langages reconnaissables sont stables par les
  constructions de la classe des langages rationnels.
\item étant donné un automate $\varAuA$, construire une expression rationnelle
  $e$ dont l'interprétation est le langage reconnu par $\varAuA$.
\end{itemize}

Pour la première étape, nous aurons besoin de faire des constructions sur les
automates.

\subsection{Constructions d'automates}

Les langages rationnels contiennent les singletons $\{\varla\}$ pour
$\varla \in \ABA$, l'ensemble vide et sont stables par trois opérations~:
\begin{itemize}
\item l'union,
\item la concaténation de langage,
\item l'étoile d'un langage.
\end{itemize}

Il nous faut donc construire des automates représentant tous ces cas. Traitons
d'abord les cas les plus simples.

Il est clair qu'on peut construire des automates pour les singletons
(\latinexpr{c.f.} \cref{fig.auto.a}) et pour l'ensemble vide
(\latinexpr{c.f.} \cref{fig.auto.vide}).

\begin{figure}[t]
  \centering
  \begin{tikzpicture}[node distance = 2cm, on grid, auto]
      \node[state, initial] (q_0) {$\varqP{0}$};
      \node[state,accepting] (q_1) [right = of q_0] {$\varqP{1}$};
      \path[->,>=latex]
      (q_0) edge [bend left] node {$\varla$} (q_1);
  \end{tikzpicture}

  \vspace{0.2cm}
  \hrule

  \vspace{0.2cm}
  \caption{Automate $\varAuA_{\varla}$ reconnaissant exactement $\varla$}
  \label{fig.auto.a}
\end{figure}

\begin{figure}[t]
  \centering
  \hrule

  \vspace{0.2cm}
  \begin{tikzpicture}[node distance = 2cm, on grid, auto]
      \node[state, initial] (q_0) {$\varqP{0}$};
  \end{tikzpicture}

  \vspace{0.2cm}
  \hrule

  \vspace{0.2cm}
  \caption{Automate $\varAuA_{\vide}$ reconnaissant $\vide$}
  \label{fig.auto.vide}
\end{figure}

Le point essentiel est, par conséquent, la stabilité par les trois opérations.
Il nous faut pour cela introduire deux constructions d'automates~: l'automate
union, et les $\Ewd$-transitions.

\begin{definition}[Automate union]
  Soient deux automates finis déterministes
  $\varAuA = (\varQA,\varqZ,\trAA,\varAAF)$ et
  $\varAuB = (\varQB,\varqZ',\trAB,\varABF)$ deux automates sur un alphabet
  $\ABA$. On définit l'automate union $\UnAuB$ par~:
  \begin{itemize}
    \begin{minipage}{0.45\textwidth}
    \item $\varQP{\UnAuB} \defeq \varQA\times\varQB$
    \item $\varqP{0,\UnAuB} \defeq (\varqZ,\varqZ')$
    \end{minipage}
    \begin{minipage}{0.45\textwidth}
    \item
      $\trAP{\UnAuB}((\varq,\varq'),\varla)\defeq
      (\trAA(\varq,\varla),\trAB(\varq',\varla))$
    \item $\varAFP{\UnAuB} \defeq \varAAF\times\varQB \cup \varQA\times \varABF$
    \end{minipage}
  \end{itemize}
\end{definition}

\begin{proposition}\label{prop.union.ratio}
  Pour tous automates $\varAuA, \varAuB$ sur un alphabet $\ABA$ donné, on a
  l'équation
  \[\varLL_{\varAuA} \cup \varLL_{\varAuB} = \varLL_{\UnAuB}\]
\end{proposition}

\begin{proof}
  On montre d'abord par récurrence sur $\varwd$ que pour tout
  $\varwd \in \toStA$, on a
  \[\trASP{\UnAuB}((\varqZ,\varqZ'),\varwd) =
  (\trAAS(\varqZ,\varwd),\trABS(\varqZ',\varwd))\]
  \begin{itemize}
  \item pour $\varwd = \Ewd$, le résultat est par définition.
  \item si l'égalité est vraie pour $\varwd$ et pour tout $\varla \in \ABA$, on
    a
    \begin{align*}
      \trASP{\UnAuB}((\varqZ,\varqZ'),\varwd\wdPt\varla) &=
      \trAP{\UnAuB}(\trASP{\UnAuB}((\varqZ,\varqZ'),\varwd),\varla)\\
      &= \trAP{\UnAuB}((\trAAS(\varqZ,\varwd),\trABS(\varqZ',\varwd)),\varla)\\
      &= (\trAA(\trAAS(\varqZ,\varwd),\varla),
      \trAB(\trABS(\varqZ',\varwd),\varla))\\
      &= (\trAAS(\varqZ,\varwd\wdPt\varla),\trABS(\varqZ',\varwd\wdPt\varla))
    \end{align*}
  \end{itemize}

  De plus, $(\varq,\varq') \in \varAAF\times \varQB \cup \varQA\times\varABF$
  signifie que $\varq \in \varAAF$ ou $\varq' \in \varABF$. On en déduit donc,
  en utilisant le résultat précédent, que
  \[\UnAuB\satisf \varwd \iff
  \varAuA \satisf \varwd \lor \varAuB\satisf \varwd\qedhere\]
\end{proof}

\begin{exercise}[Automate produit]
  Soient $\varAuA$ et $\varAuB$ deux automates sur un alphabet $\ABA$. On
  définit l'automate produit $\varAuA \times \varAuB$ de la même façon que
  pour l'automate $\UnAuB$ à l'exception de l'ensemble des états terminaux,
  qui est donné par
  \[\varAFP{\varAuA\times\varAuB} \defeq \varAAF \times \varABF\]
  Montrer que cet automate reconnait $\varLL_{\varAuA} \cap \varLL_{\varAuB}$, en
  déduire que les langages reconnaissables sont stables par intersection
  (binaire).

  En s'inspirant de la construction précédente, montrer que pour toute opération
  booléenne $\makeFunName{f}{\intB^2}{\intB}$ et tous langages $\varLL,\varLL'$
  reconnaissables par automates, le langage
  \[f(\varLL,\varLL') \defeq
  \compr{\varwd \in \toStA}
        {f(\indicP{\varLL}(\varwd),\indicP{\varLL'}(\varwd)) = 1}\]
  est aussi un langage reconnaissable.
\end{exercise}

On veut maintenant montrer la stabilité par concaténation de deux langages et
la stabilité par étoile d'un langage. L'idée derrière ces deux constructions est
relativement simple~: pour deux automates $\varAuA,\varAuB$, l'automate
reconnaissant la concaténation est une copie de chacun des deux automates, en
ajoutant des transitions entre les états acceptants de $\varAuA$ et l'état
initial de $\varAuB$. De même pour l'étoile d'un langage, on ajoute un nouvel
état initial, acceptant, on le relie à l'ancien état initial et on crée une
transition entre chaque état acceptant et le nouvel état initial.

Cependant, ces transitions peuvent être effectuées automatiquement, en ayant la
possibilité pour un état de faire (ou non) cette transition avant la lecture
d'une nouvelle lettre. Cela nous pousse à ajouter deux choses~: le non
déterminisme, pour pouvoir choisir ou non si l'on fait une transition, et les
$\Ewd$-transitions, qui sont des transitions qui, comme décrites précédemment,
peuvent s'exécuter ou non et sans lecture de lettre.

Commençons par présenter un automate non déterministe~: celui-ci est un automate
dans lequel, au lieu d'une fonction de transition indiquant à partir d'un état
et d'une lettre quel est l'état lu ensuite, on dispose d'une relation de
transition indiquant pour deux états et une lettre s'il existe une transition
entre ces deux états étiquetée par cette lettre. On lui préfère en général une
présentation équivalente~: à un état et une lettre on associe l'ensemble des
états possibles à atteindre.

\begin{definition}[Automate fini non déterministe]
  Soit $\ABA$ un alphabet. On appelle automate fini non déterministe $\varAuA$
  un quadruplet $(\varQQ,\varQZ,\trA,\varAF)$ où~:
  \begin{itemize}
  \item $\varQQ$ est un ensemble fini appelé ensemble des états de $\varAuA$,
  \item $\varQZ\subseteq\varQQ$ est l'ensemble des états initiaux,
  \item $\makeFunName{\trA}{\varQQ\times\ABA}{\powerset(\varQQ)}$ est la
    fonction de transition non déterministe,
  \item $\varAF\subseteq\varQQ$ est l'ensemble des états terminaux.
  \end{itemize}
\end{definition}

\begin{remark}
  Dans le cas d'un automate non déterministe, la partialité de la fonction de
  transition se traduit par le fait que $\trA(\varq) = \vide$, ce qui rend la
  fonction (de codomaine $\powerset(\varQQ)$) totale.
\end{remark}

On peut définir la fonction de transition étendue, dont le principe est de
considérer l'union à chaque fois qu'on prend un ensemble d'états possibles.

\begin{definition}[Fonction de transition étendue]
  Soit un automate fini non déterministe $\varAuA$ sur un alphabet $\ABA$,
  on note $\makeFunName{\trAS}{\varQQ \times \toStA}{\powerset(\varQQ)}$ sa
  fonction de transition étendue, définie par induction par~:
  \begin{itemize}
  \item $\trAS(\varq,\Ewd) = \{\varq\}$
  \item $\displaystyle\trAS(\varq,\varwd\wdPt\varla) =
    \bigcup_{\varq' \in \trAS(\varq,\varwd)} \trA(\varq',\varla)$
  \end{itemize}
\end{definition}

\begin{definition}[Langage accepté]
  Soit $\ABA$ un alphabet et $\varAuA$ un automate fini non déterministe sur
  $\ABA$. On dit que $\varAuA$ accepte (ou reconnaît) le mot $\varwd\in\toStA$
  s'il existe $\varqZ \in \varQZ$ tel que
  $\trAS(\varqZ,\varwd) \cap \varAF\neq\vide$, ce qu'on note
  $\varAuA \satisf \varwd$.
  On définit le langage reconnu par $\varAuA$ par~:
  \[\varLL_{\varAuA}\defeq \compr{\varwd \in \toStA}{\varAuA \satisf\varwd}\]
\end{definition}

Un automate fini non déterministe se représente comme un automate fini
déterministe, à ceci près qu'on accepte d'avoir plusieurs flèches depuis un
état qui sont étiquetées par la même lettre, et qu'on peut avoir plusieurs
états initiaux. Accepter un mot dans un tel automate est alors simplement
trouver au moins un chemin partant d'un état initial et arrivant à un état
terminal en suivant des transitions étiquetées par les lettres du mot lu.

Cela mène à une définition alternative de l'acceptation d'un mot, qu'on traite
en exercice. La notion de trace est assez générale, et représente l'historique
de l'exécution d'une entrée sur une machine~: ici, la trace d'un automate est
alors la suite des états dans lesquels est la machine en lisant un mot.

\begin{exercise}\label{exo.traces.auto}
  Soit $\ABA$ un alphabet et $\varAuA$ un automate fini non déterministe sur
  $\ABA$. On définit, pour tout $\varwd \in \toStA$, l'ensemble $\tr(\varwd)$
  des traces de $\varwd$ sur $\varAuA$, comme l'ensemble des suites finies
  $(\varqP{0},\ldots,\varqP{n})$ d'états de $\varAuA$ dont le premier élément
  est un élément de $\varQZ$ et tel que $\trA(\varqP{i},\varwd_i) = \varqP{i+1}$
  pour tout $i \in \intervInt{n}$. Montrer que $\varAuA \satisf\varwd$ si et
  seulement s'il existe un élément $p \in \tr(\varwd)$ dont le dernier état est
  un état terminal.
\end{exercise}

On pourrait s'attendre à ce que relâcher tant de contraintes augmente la
puissance des automates, mais ce n'est en réalité pas le cas. Il est clair qu'un
automate déterministe peut être transformé en un automate non déterministe, en
prenant $\varQZ \defeq \{\varqZ\}$ et pour fonction de transition non
déterministe le singleton de l'image par la fonction de transition déterministe.
La réciproque, elle, demande une construction plus subtile, appelée l'automate
des parties.

\begin{definition}[Automate des parties]
  Soit un alphabet $\ABA$ et un automate fini non déterministe $\varAuA$ sur
  $\ABA$. On définit l'automate des parties de $\varAuA$, noté
  $\powerset(\varAuA)$, par~:
  \begin{itemize}
    \begin{minipage}{0.45\textwidth}
    \item $\varQP{\powerset(\varAuA)} \defeq \powerset(\varQP{\varAuA})$
    \item $\varqP{0,\powerset(\varAuA)} \defeq \varQP{0,\varAuA}$
    \end{minipage}
    \begin{minipage}{0.45\textwidth}
      \item $\trAP{\powerset(\varAuA)} (P,\varla) \defeq
        \bigcup_{\varq \in P} \trAP{\varAuA}(\varq,\varla)$
      \item $\varAFP{\powerset(\varAuA)} \defeq
        \compr{P \subseteq \varQP{\varAuA}}{P\cap \varAFP{\varAuA} \neq \vide}$
    \end{minipage}
  \end{itemize}
\end{definition}

\begin{proposition}
  Soit un automate non déterministe $\varAuA$ sur un alphabet $\ABA$, on a
  l'égalité
  \[\varLL_{\varAuA} = \varLL_{\powerset(\varAuA)}\]
\end{proposition}

\begin{proof}
  On prouve par induction sur $\varwd$ que pour tout $\varwd \in \toStA$, on a
  l'égalité suivante pour tout $P \subseteq \varQP{\varAuA}$~:
  \[\trASP{\powerset(\varAuA)}(P,\varwd) =
  \bigcup_{\varq \in P}\trAAS(q,\varwd)\]
  (où le deuxième terme est une image direct).
  \begin{itemize}
  \item dans le cas où $\varwd = \Ewd$, le résultat est direct.
  \item supposons que l'hypothèse d'induction est valide pour $\varwd$ et
    montrons le cas $\varwd\wdPt\varla$~:
    \begin{align*}
      \trASP{\powerset(\varAuA)}(P,\varwd\wdPt\varla) &=
      \bigcup_{\varq \in \trASP{\powerset(\varAuA)}(P,\varwd)}
      \trAP{\varAuA}(\varq,\varla)\\
      &= \bigcup_{\varq \in \bigcup_{\varq' \in P}\trAAS(\varq',\varwd)}
      \trAA(\varq,\varla)\\
      &= \bigcup_{\varq' \in P}\bigcup_{\varq \in \trAAS(\varq',\varwd)}
      \trAA(\varq,\varla)\\
      &= \bigcup_{\varq \in P}\trAAS(\varq,\varwd\wdPt\varla)
    \end{align*}
  \end{itemize}
  D'où le résultat par induction. On voit alors que la condition pour que
  $\varwd$ soit accepté, dans les deux automates, revient à
  \[\trASP{\powerset(\varAuA)} (\varQZ,\varwd) \cap\varAAF \neq \vide\]
  donc $\varLL_{\powerset(\varAuA)} = \varLL_{\varAuA}$.
\end{proof}

\begin{remark}
  Cette construction n'est pas gratuite en pratique~: elle fait passer d'un
  ensemble d'états à son ensemble des parties, ce qui donne une complexité
  exponentielle à la construction.
\end{remark}

On peut donc utiliser des automates non déterministes pour étudier les langages
reconnaissables. On introduit maintenant les $\Ewd$-transitions, qui
ont plus de sens dans le contexte des automates non déterministes.

\begin{definition}[Automate avec $\Ewd$-transitions]
  On définit un automate fini non déterministe avec $\Ewd$-transitions comme un
  automate fini non déterministe, mais dans lequel la fonction de transition est
  de la forme
  $\makeFunName{\trA}{\varQQ\times(\ABA \cup \{\Ewd\})}{\powerset(\varQQ)}$.
\end{definition}

Pour définir l'acceptation d'un mot par un automate avec $\Ewd$-transitions, on
utilisera le formalisme des traces, introduit dans l'\cref{exo.traces.auto}.

\begin{definition}[Langage reconnu par un automate avec $\Ewd$-transitions]
  Soit $\ABA$ un alphabet et $\varAuA$ un automate non déterministe avec
  $\Ewd$-transitions sur $\ABA$. On définit l'ensemble $\tr(\varwd)$ des
  traces sur un mot $\varwd$ comme l'ensemble des suites finies
  $(\varqZ,\ldots,\varqP{n})$ d'états de $\varQQ$ où $\varqZ \in \varQZ$ et
  telles qu'il existe pour tout $i \in \intervInt{n}$ une lettre (possiblement
  $\Ewd$) $\varla_i \in \ABA\cup\{\Ewd\}$ telle que les deux conditions
  suivantes sont vérifiées~:
  \begin{itemize}
  \item $\varqP{i+1} \in \trA(\varqP{i},\varla_i)$
  \item $\varwd = \varla_0\wdPt\cdots\wdPt \varla_{n-1}$
  \end{itemize}
  On dit que $\varAuA$ accepte le mot $\varwd$ s'il existe une suite
  $(\varqZ,\ldots,\varqP{n}) \in \tr(\varwd)$ telle que $\varqP{n} \in \varAF$.

  On définit le langage reconnu par $\varAuA$ comme l'ensemble des mots reconnus
  par $\varAuA$.
\end{definition}

Les traces dans un automate avec $\Ewd$-transitions représentent simplement les
exécutions possibles lors de la lecture d'un mot, mais où il est possible
d'effectuer gratuitement des transitions étiquetées par $\Ewd$.

Là encore, il est clair qu'un automate sans $\Ewd$-transition peut être simulé
par un automate avec $\Ewd$-transition, en n'ajoutant aucune $\Ewd$-transition.

On introduit maintenant l'élimination des $\Ewd$-transitions, permettant de
construire un automate sans $\Ewd$-transition mais reconnaissant le même
langage.

L'idée de l'élimination des $\Ewd$-transitions est assez naturelle~: au lieu
de faire des sauts spontanés dans une $\Ewd$-transition, on décide de faire une
grande transition contenant toutes les $\Ewd$-transitions ainsi que la
transition qui suit. Il nous faut aussi modifier l'ensemble des états terminaux
pour ajouter les états qui peuvent atteindre un état terminal uniquement en
effectuant des $\Ewd$-transitions.

\begin{definition}[\'Elimination des $\Ewd$-transitions]
  Soit $\ABA$ un alphabet et $\varAuA$ un automate avec $\Ewd$-transitions sur
  $\ABA$. On définit la relation $\rtClot{\to}$ sur les états de $\varAuA$ comme
  la plus petite relation  réflexive et transitive contenant la relation
  \[\varq\to \varq' \defeq \varq' \in \trA(\varq,\Ewd)\]

  On définit l'automate $\varAuP{\nEwd}$ non déterministe sans $\Ewd$-transition
  par la construction suivante~:
  \begin{itemize}
  \item $\varQP{\varAuP{\nEwd}} \defeq \varQP{\varAuA}$
  \item $\varQP{0,\varAuP{\nEwd}} \defeq \varQP{0,\varAuA}$
  \item $\displaystyle\trAP{\varAuP{\nEwd}}(\varq,\varla) \defeq
    \bigcup_{\genfrac{}{}{0pt}{1}{\varq' \in \varQP{\varAuA}}{\varq\rtClot{\to}\varq'}}
    \trAP{\varAuA}(\varq',\varla)$
  \item $\varAFP{\varAuP{\nEwd}} =
    \compr{\varq \in \varQQ}
          {\exists \varq'\in\varQQ, \varq \rtClot{\to}\varq' \land
            \varq' \in \varAFP{\varAuA}}$
  \end{itemize}
\end{definition}

\begin{proposition}\label{prop.epsi.equiv}
  Soit $\ABA$ un alphabet et $\varAuA$ un automate avec $\Ewd$-transitions sur
  $\ABA$. Alors
  \[\varLL_{\varAuA} = \varLL_{\varAuP{\nEwd}}\]
\end{proposition}

\begin{proof}
  On prouve d'abord que $\varq'\in \trAP{\varAuP{\nEwd}}(\varq,\varla)$ si et
  seulement s'il existe une suite finie $(\varqZ,\ldots,\varqP{n})$ telle que~:
  \begin{itemize}
  \item $\varqZ = \varq$
  \item $\forall i \in \{0,\ldots,n-2\}, \varqP{i+1} \in \trAA(\varqP{i},\Ewd)$
  \item $\varqP{n} \in \trAA(\varqP{n-1},\varla)$
  \item $\varqP{n} = \varq'$
  \end{itemize}
  Pour un sens, on remarque que dans une telle suite,
  $\varqZ \rtClot{\to} \varqP{n-1}$, ce qui nous donne l'inclusion
  $\trAP{\varAuP{\nEwd}}(\varq,\varla)\supseteq \trAA(\varqP{n-1},\varla)$.
  Pour l'autre sens, on sait par définition qu'il existe $\varq''$ tel que
  $\varq\rtClot{\to} \varq''$ et $\varq'\in\trAA (\varq'',\varla)$, donc il nous
  reste à vérifier que si $\varq\rtClot{\to}\varq''$ alors il existe une suite
  d'états uniquement reliés par des $\Ewd$-transition, ce qui se fait par
  induction sur $\rtClot{\to}$ (on laisse la vérification au lecteur).

  Ainsi, on a une trace $(\varqZ,\ldots,\varqP{n})$ dans $\varAuA$ pour un mot
  $\varwd$ donné si et seulement si $(\varqZ',\ldots,\varqP{p}')$, obtenue en
  supprimant les états reliés par des $\Ewd$-transitions, est une trace pour
  $\varwd$ dans $\varAuP{\nEwd}$ (réciproquement si on peut ajouter des états
  reliés par des $\Ewd$-transitions).

  De plus, un état $\varqP{n}$ est terminal dans $\varAuP{\nEwd}$ si et
  seulement s'il existe une suite d'$\Ewd$-transitions menant à un état
  terminal~: ainsi une trace terminant sur un tel état dans $\varAuA$ donne une
  trace acceptante en ajoutant des $\Ewd$-transitions. Réciproquement, on
  récupère une trace acceptante dans $\varAuP{\nEwd}$ en supprimant ces
  $\Ewd$-transitions. Ainsi,
  \[\varLL_{\varAuA} = \varLL_{\varAuP{\nEwd}}\qedhere\]
\end{proof}

On sait donc que les automates déterministes, non déterministes et
non déterministes avec $\Ewd$-transitions, sont tous équivalents du
point de vue de la classe de langage qu'ils reconnaissent.

On peut maintenant prouver la stabilité des langages reconnaissables par
concaténation de langages et par étoile.

\begin{proposition}\label{prop.concat.ratio}
  Soit $\ABA$ un alphabet et $\varAuA, \varAuB$ deux automates non déterministes
  avec $\Ewd$-transitions sur $\ABA$. Alors il existe un automate reconnaissant
  \[\varLL_{\varAuA} \wdPt \varLL_{\varAuB}\]
\end{proposition}

\begin{proof}
  On construit l'automate $\varAuA''$ dont l'ensemble des états est la somme
  directe des états de $\varAuA$ et $\varAuB$ (pour simplifier les notations on
  va supposer les deux ensembles d'états disjoints et prendre leur union), les
  transitions sont uniquement celles définies au sein de $\varAuA$ et $\varAuB$
  sauf entre chaque état $\varqP{f}$ de $\varAuA$ et chaque état initial
  $\varqP{i}$ de $\varAuB$ où l'on rajoute une transition
  $\varqP{f}\xrightarrow{\Ewd}\varqP{i}$. Les états initiaux sont ceux de
  $\varAuA$ et les états terminaux ceux de $\varAuB$.

  Accepter un mot dans cet automate signifie donc avoir une trace
  $(\varqZ,\ldots,\varqP{n})$ commençant en un état initial de $\varAuA$ et
  terminant dans un état terminal de $\varAuB$. Comme la seule façon de passer
  d'un état de $\varAuA$ a un état de $\varAuB$ est d'utiliser une des
  $\Ewd$-transitions ajoutées, on sait que la trace peut se décomposer en
  $(\varqZ,\ldots,\varqP{p})$ et $(\varqP{{p+1}},\ldots,\varqP{n})$ et ces deux
  traces sont acceptantes respectivement pour un mot dans $\varLL_{\varAuA}$ et
  pour un mot dans $\varLL_{\varAuB}$, donc la trace est acceptant pour un mot
  dans $\varLL_{\varAuA} \wdPt \varLL_{\varAuB}$.

  Réciproquement, si on a deux traces $(\varqZ,\ldots,\varqP{p})$ et
  $(\varqP{p+1},\ldots,\varqP{n})$ pour respectivement
  $\varwd \in \varLL_{\varAuA}$ et $\varwv\in\varLL_{\varAuB}$, alors
  $(\varqZ,\ldots,\varqP{n})$ est une trace acceptante pour $\varwd\wdPt\varwv$.
\end{proof}

On prouve de même qu'ajouter simplement des transitions entre les états
terminaux et initiaux d'un automate donne le l'étoile du langage. Comme la
preuve est très proche de la preuve précédente, nous préférons la donner en
exercice.

\begin{definition}[Automate étoile]
  Soit $\ABA$ un alphabet et $\varAuA$ un automate non déterministe sur $\ABA$.
  On définit l'automate étoile $\toSt{\varAuA}$ par~:
  \begin{itemize}
  \item $\varQP{\toSt{\varAuA}} \defeq \varQA \sqcup \{p\}$
  \item $\varQP{0,\toSt{\varAuA}} \defeq \{p\}$
  \item $\trAP{\toSt{\varAuA}}$ est définie comme $\trAA$ mais en ajoutant des
    $\Ewd$-transitions entre $p$ et chaque $\varqZ\in\varQP{0,\varAuA}$, ainsi
    qu'entre chaque $\varqP{f} \in \varAAF$ et $p$.
  \item $\varAFP{\toSt{\varAuA}} \defeq \{p\}$.
  \end{itemize}
\end{definition}

\begin{exercise}\label{prop.etoile.ratio}
  Soit $\ABA$ un alphabet et $\varAuA$ un automate sur $\ABA$, montrer que
  \[\varLL_{\toSt{\varAuA}} = \toSt{\varLL_{\varAuA}}\]
\end{exercise}

\noindent
On déduit donc des constructions effectuées la stabilité par les opérations
rationnelles.

\begin{corollary}
  Soit $\ABA$ un alphabet. La classe $\Reco(\ABA)$ contient la classe
  $\Ratio(\ABA)$.
\end{corollary}

\begin{proof}
  On a déjà dit que $\vide \in \Reco(\ABA)$ et que $\{\varla\}\in \Reco(\ABA)$
  pour tout $\varla \in \ABA$. De plus, en utilisant la \cref{prop.union.ratio},
  la \cref{prop.concat.ratio} et l'\cref{prop.etoile.ratio}, on en déduit que
  $\Reco(\ABA)$ est stable par union, concaténation et étoile. On en déduit
  donc, puisque $\Ratio(\ABA)$ est la plus petite telle classe, que
  $\Ratio(\ABA)\subseteq\Reco(\ABA)$.
\end{proof}

Avant de montrer l'inclusion réciproque, il nous faut nous attarder sur les
propriétés des automates.

\subsection{Propriétés d'automates}

Plusieurs propriétés d'intérêt existent, à propos des automates. Pour faciliter
la génération d'une expression rationnelle à pratir d'un automate, nous allons
présenter des propriétés basiques à leur propos. La première est celle
d'avoir une table de transition totale plutôt que partielle.

\begin{definition}[Automate complet]
  Un automate $\varAuA$ est dit complet si la fonction de transition $\trA$ est
  une fonction totale.
\end{definition}

Il est facile d'assurer cette propriété en ajoutant un état dit \og puits\fg
vers lequel pointent toutes les transitions à ajouter, et les transitions
partant de cet état.

\begin{proposition}
  Soit $\varAuA$ un automate sur un alphabet $\ABA$, il existe un automate
  $\varAuB$ qui est complet et reconnaît le même langage que $\varAuA$.
\end{proposition}

\begin{proof}
  On construit l'automate $\varAuB$ comme décrit ci-dessus~:
  \begin{itemize}
    \begin{minipage}{0.3\textwidth}
    \item $\varQQ' \defeq \varQQ \sqcup \{p\}$
    \end{minipage}
    \begin{minipage}{0.3\textwidth}
    \item $\varqZ' \defeq \varqZ$
    \end{minipage}
    \begin{minipage}{0.3\textwidth}
    \item $\varAF' \defeq \varAF$
    \end{minipage}
  \item $\makeFunName{\trA'}{\varQQ'\times\ABA}{\varQQ'}$ est définie par
    disjonction de cas~:
    \begin{itemize}
    \item si $(\varq,\varla) \in \dom(\trA)$, alors
      $\trA'(\varq,\varla) \defeq \trA(\varq,\varla)$
    \item sinon, $\trA'(\varq,\varla) \defeq p$
    \end{itemize}
  \end{itemize}

  On veut maintenant prouver que les deux automates reconnaissent le même
  langage. Soit $\varwd \in \toStA$, montrons que $\trAS(\varqZ,\varwd)$ est
  indéfini si et seulement si $\toSt{\trA'}(\varqZ,\varwd) = p$ et que, sinon,
  $\toSt{\trA'}(\varqZ,\varwd) = \trAS(\varqZ,\varwd)$, par induction sur
  $\varwd$~:
  \begin{itemize}
  \item dans le cas de $\Ewd$, on sait que
    $\toSt{\trA'}(\varqZ,\Ewd) = \trA(\varqZ,\Ewd) = \varqZ$.
  \item supposons l'hypothèse vraie pour $\varwd$, et soit $\varla \in \ABA$.
    
    Montrons d'abord que si $\trAS(\varqZ,\varwd)$ est indéfinie, alors
    $\toSt{\trA'}(\varqZ,\varwd) = p$.
    La valeur $\trAS(\varqZ,\varwd)$ peut être définie ou indéfinie, nous
    donnant deux cas~:
    \begin{itemize}
    \item Si $\trAS(\varqZ,\varwd)$ est indéfinie, alors par hypothèse
      d'induction $\toSt{\trA'}(\varqZ,\varwd) = p$ et comme
      $\trA'(p,\varla) = p$, on en déduit que
      $\toSt{\trA'}(\varqZ,\varwd\wdPt\varla) = p$.
    \item Si $\trAS(\varqZ,\varwd) = \varq$ et $\trA(\varq,\varla)$ est
      indéfini, alors  $\toSt{\trA'}(\varqZ,\varwd) = \varq$ par hypothèse
      d'induction puis, par définition, $\toSt{\trA'}(\varq,\varla) = p$.
    \end{itemize}
    
    Réciproquement, si $\toSt{\trA'}(\varqZ,\varwd\wdPt\varla) = p$ alors c'est
    que la transition a atteint un état puits à un moment, et cela correspondait
    à une transition non définie dans $\trA$ (la preuve est parfaitement
    analogue à celle du sens direct).

    Dans l'autre cas, si $\toSt{\trA'}(\varqZ,\varwd\wdPt\varla) = \varq$ où
    $\varq \in \varQQ$, alors par hypothèse d'induction on trouve $\varq'$ tel
    que $\toSt{\trA'}(\varqZ,\varwd) = \trA(\varqZ,\varwd) = \varq'$, donc
    $\trA'(\varq',\varla) = \trA(\varq',\varla)$ par définition, d'où
    $\toSt{\trA'}(\varqZ,\varwd\wdPt\varla)=\trAS(\varqZ,\varwd\wdPt\varla)$.
  \end{itemize}

  On en déduit donc~:
  \[\trAS(\varqZ,\varwd) \in \varAF \iff \toSt{\trA'}(\varqZ,\varwd)\in \varAF
  \qedhere\]
\end{proof}

Une autre propriété importante sur un automate est celle d'être émondé.
Pour donner une idée de leur intérêt, reprenons la construction précédente.
En complétant un automate, on ajoute un état puits, mais arriver à cet état
signifie forcément que la lecture du mot n'arrivera jamais à un état acceptant.
On peut donc, si on préfère minimiser le nombre d'états qui apparaissent
effectivement dans l'automate, décider de supprimer notre état puits, mais aussi
tous les états qui n'arriveront jamais à un état acceptant. De tels états sont
dits co-accessibles. De façon duale, un état accessible est un état qu'on peut
atteindre à la lecture d'un mot et ils sont eux aussi les seuls d'intérêt~:
en effet, même si un état mène à un état acceptant, cela n'a pas d'importance
si aucun chemin ne mène à lui.

On peut donc donner des définitions similaires à propos d'un automate pour
parler de tous ses états~: un automate accessible est un automate dont tous les
états sont accessibles et un automate co-accessible est un automate dont tous
les états sont co-accessibles. Un automate à la fois accessible et co-accessible
est appelé un automate émondé, et il est toujours possible de construire un
automate émondé, en ne gardant que les états traversés lors de traces
acceptantes.

\begin{proposition}
  Soit $\ABA$ un alphabet et $\varAuA$ un automate sur $\ABA$. Alors il existe
  un automate $\varAuB$ émondé qui reconnaît le même langage que $\varAuA$.
\end{proposition}

\begin{proof}
  On remarque que pour toute trace $(\varqZ,\ldots,\varqP{n})$ acceptante, les
  états $\varqP{i}, i \in \{0,\ldots,n\}$ sont à la fois accessibles et
  co-accessibles. Ainsi, si une trace contient un état non accessible ou non
  co-accessible, alors elle n'est pas acceptante. Puisqu'un mot $\varwd$ n'est
  accepté que lorsqu'il existe une trace acceptante pour $\varwd$, l'ensemble
  des mots acceptés n'est pas changé en ne gardant que les états à la fois
  accessibles et co-accessibles. On en déduit que l'automate $\varAuB$ induit
  par $\varAuA$ en ne prenant que les états intervenant dans des traces
  acceptantes reconnaît le même langage que $\varAuA$.
\end{proof}

Enfin, parlons de la notion d'automate normalisé. Un tel automate est un
automate dans lequel on peut trouver trois parties~: l'état initial,
unique, le c\oe ur de l'automate et l'état terminal. Celui-ci permet d'avoir
une forme plus linéaire à un automate.

\begin{definition}[Automate normalisé]
  Soit $\ABA$ un alphabet. Un automate normalisé sur $\ABA$ est un automate sur
  $\ABA$ tel qu'il existe un unique état initial, un unique état terminal, tel
  qu'il n'existe aucune transition d'un état vers l'état initial et aucune
  transition depuis l'état terminal.
\end{definition}

\begin{proposition}
  Soit $\ABA$ un alphabet et $\varAuA$ un automate sur $\ABA$, il existe alors
  un automate non déterministe avec $\Ewd$-transitions, normalisé, reconnaissant
  le même langage que $\varAuA$.
\end{proposition}

\begin{proof}
  On construit l'automate simplement en ajoutant un nouvel état initial $\varqZ$
  et un nouvel état terminal $\varqP{f}$ et des $\Ewd$-transitions depuis
  $\varqZ$ vers tout état anciennement initial, et de tout état anciennement
  terminal vers $\varqP{f}$.
\end{proof}

\begin{exercise}
  Montrer que l'automate construit précédemment reconnaît le même langage que
  $\varAuA$.
\end{exercise}

On peut aussi procéder à l'élimination des $\Ewd$-transitions sur un tel
automate, mais on remarque alors un problème~: si le langage contient $\Ewd$,
alors il doit y avoir une transition de $\varqZ$ à $\varqP{f}$ étiquetée par
$\Ewd$, ce qui ne peut pas s'éliminer en gardant $\varqZ$ et $\varqP{f}$
distincts. On a donc besoin d'ajouter que le langage reconnu par $\varAuA$ ne
contient par $\Ewd$, si l'on veut renforcer le résultat en donnant un automate
sans $\Ewd$-transitions normalisé.

On montre maintenant que $\Reco(\ABA)\subseteq \Ratio(\ABA)$. Pour montrer cette
inclusion, on va construire une expression rationnelle par un processus
d'élimination d'états depuis un automate. \'Etant donné un automate $\varAuA$,
on peut considérer les étiquettes comme des expressions rationnelles limitées à
un caractère. Avec cette considération, effectuer deux transitions
$\xrightarrow a \xrightarrow b$ correspond à une seule transition
$\xrightarrow{ab}$. Cela nous donne alors une procédure pour regrouper des
transitions, et ainsi éliminer des états (puisque l'état du milieu entre la
transition $a$ et la transition $b$ peut alors être supprimé). On a ainsi un
algorithme qui construit un expression rationnelle pour un automate. On remarque
que les $\Ewd$-transitions sont parfaitement adaptées dans ce contexte, puisque
$\Val(\Ewd e) = \Val(e)$, comme on l'attend d'une $\Ewd$-transition.

La procédure à suivre est donc la suivante~:
\begin{itemize}
\item définir une notion plus générale d'automate fini qui permet d'utiliser des
  expressions rationnelles pour les étiquettes,
\item créer le processus d'élimination d'état,
\item montrer que ce processus conserve le langage reconnu.
\end{itemize}

\begin{definition}[Automate fini généralisé]
  Soit $\ABA$ un alphabet, on appelle un automate fini généralisé un
  quadruplet $(\varQQ,\varQZ,\trA,\varAF)$ tel que~:
  \begin{itemize}
  \item $varQQ$ est l'ensemble fini des états,
  \item $\varQZ\subseteq\varQQ$ est l'ensemble des états initiaux,
  \item $\makeFunName{\trA}{\varQQ \times \Regex(\ABA)}{\powerset(\varQQ)}$
    est la fonction de transition généralisée,
  \item $\varAF\subseteq\varQQ$ est l'ensemble des états terminaux.
  \end{itemize}
\end{definition}

\begin{definition}[Acceptation d'un mot par un automate généralisé]
  Soit $\ABA$ un alphabet et $\varAuA$ un automate généralisé sur $\ABA$. On dit
  qu'un mot $\varwd\in \toStA$ est accepté (ou reconnu) par $\varAuA$ s'il
  existe une décomposition de $\varwd$ en une suite de mots
  $\varwd = \varwd_0\wdPt\cdots\wdPt \varwd_{p-1}$ et une suite d'états
  $(\varqZ,\ldots,\varqP{p})$ telles que~:
  \begin{itemize}
    \begin{minipage}{0.45\textwidth}
    \item $\varqZ \in \varQZ$
    \end{minipage}
    \begin{minipage}{0.45\textwidth}
    \item $\varqP{p} \in \varAF$
    \end{minipage}
  \item $\forall i \in \intervInt{p}, \exists e \in \Regex(\ABA),
    \varwd_i \in \Val(e) \land \varqP{i+1}\in \trA(\varqP{i},e)$
  \end{itemize}
\end{definition}

On voit ici qu'une inclusion évidente existe entre les automates non
déterministes et les automates généralisés~: il suffit de transformer
l'étiquette $\varla$ par le regex $\varla$. On veut donc maintenant construire
le processus d'élimination. On a dit précédemment qu'une suite de deux
transitions peut se transformer en une seule concaténation, ce qui donne l'idée
suivante de transformation~:

\begin{figure}[H]
  \centering
  \begin{tikzpicture}[node distance = 2cm, on grid, auto]
    \node[state] (q_0) {$\varqZ$};
    \node[state] (q_1) [right = of q_0] {$\varqP{1}$};
    \node[state] (q_2) [right = of q_1] {$\varqP{2}$};
    \node (use) [right = of q_2] {$\longmapsto$};
    \node[state] (q_3) [right = of use] {$\varqZ$};
    \node (use2) [right = of q_3] {};
    \node[state] (q_4) [right = of use2] {$\varqP{2}$};
    \path[->,>=latex]
    (q_0) edge node {$\varla$} (q_1)
    (q_1) edge node {$\varlb$} (q_2)
    (q_3) edge node {$\varla\varlb$} (q_4);
  \end{tikzpicture}
\end{figure}

En effectuant cette transformation entre chaque paire de transitions passant par
$\varqP{1}$, il semble donc possible d'éliminer l'état $\varqP{1}$. Cependant,
un cas important reste possible~: on peut avoir une boucle sur $\varqP{1}$. Dans
ce cas, il est nécessaire d'ajouter pour chaque expression $e$ étiquetant une
boucle, l'expression $\varla \toSt{e}\varlb$ au lieu de simplement
$\varla\varlb$, comme on peut le voir dans le diagramme suivant~:

\begin{figure}[H]
  \centering
  \begin{tikzpicture}[node distance = 2cm, on grid, auto]
    \node[state] (q_0) {$q_0$};
    \node[state] (q_1) [right = of q_0] {$q_1$};
    \node[state] (q_2) [right = of q_1] {$q_2$};
    \node (use) [right = of q_2] {$\longmapsto$};
    \node[state] (q_3) [right = of use] {$q_0$};
    \node (use2) [right = of q_3] {};
    \node[state] (q_4) [right = of use2] {$q_2$};
    \path[->,>=latex]
    (q_0) edge node {$a$} (q_1)
    (q_1) edge node {$b$} (q_2)
    (q_1) edge [loop above] node {$e$} ()
    (q_3) edge node {$ae^\star b$} (q_4);
  \end{tikzpicture}
\end{figure}

Il reste un dernier souci~: il est possible qu'il existe plusieurs transitions
distinctes entre deux états donnés. Pour contrer ça, on s'assure d'abord qu'on
peut, sans perte de généralité, considérer qu'un automate généralisé n'a qu'au
plus une transition d'un état à un autre.

On introduit une nouvelle notation~: on décrit la relation de transition avec
$\varRR_{\trA}$ définie par la formule suivante~:
\[\varRR_{\trA}(\varq,e,\varq') \defeq \varq' \in \trA(\varq,e)\]
On définit aussi la notation $\sum e_i$ pour des expressions rationnelles
$e_1,\ldots,e_n$, comme $\vide$ si $n = 0$ et de façon naturelle en
itérant l'opération $+$ sinon.

\begin{proposition}
  Soit $\ABA$ un alphabet et $\varAuA$ un automate généralisé sur $\ABA$. Il
  existe alors un automate généralisé $\varAuB$, équivalent à $\varAuA$, tel que
  pour tous états $\varq,\varq'$ il existe exactement une expression $e$ telle
  que $\varRR_{\trA}(\varq,e,\varq')$.
\end{proposition}

\begin{proof}
  On construit $\varAuB$ en remplaçant pour tous états $\varq,\varq'$ les
  expressions les reliant par l'unique expressions
  \[\sum_{\genfrac{}{}{0pt}{1}{e \in \Regex(\ABA)}{\varRR_{\trA}(\varq,e,\varq')}} e\]
  Soit un mot $\varwd\in \toStA$. On considère une trace
  $(\varqZ,\ldots,\varqP{n})$ pour $\varwd$, une décomposition
  $\varwd = \varwd_0\wdPt\cdots\wdPt \varwd_{n-1}$ et une suite
  $e_0,\ldots, e_{n-1}$ telles que $\varwd_i \in \Val(e_i)$ et
  $\varRR_{\trA}(\varqP{i},e_i,\varqP{i+1})$, pour tous $i \in \intervInt{n}$.
  Cette décomposition est aussi valide dans $\varAuB$, puisque si
  $\varwd_i \in \Val(e_i)$, $u_i$ appartient à la somme des $e$ pour
  $\varRR_{\trA}(\varqP{i},e,\varqP{i+1})$. Réciproquement, appartenir à la somme
  signifie qu'on peut trouver $e$ tel que
  $\varRR_{\trA}(\varqP{i},e,\varqP{i+1})$ telle que $\varwd_i \in \Val(e_i)$. On
  en déduit donc que les traces acceptantes sont identiques dans les deux
  automates, et comme les états initiaux et terminaux sont les mêmes, $\varAuA$
  et $\varAuB$ sont équivalents.
\end{proof}

\begin{remark}
  On ajoute potentiellement de nombreuses transitions en considérant des
  transitions de la forme $\varq\xrightarrow\vide\varq'$, mais celles-ci ont le
  même rôle que l'absence de transition (puisque $\varwd \in \Val(\vide)$ n'est
  pas possible). On peut donc maintenant définir l'unique expression reliant
  deux états, ce qu'on écrit $e(\varq,\varq')$. Ainsi accepter un mot $\varwd$
  signifie obtenir une trace $(\varqZ,\ldots,\varqP{n})$ où $\varqZ$ est
  initial, $\varqP{n}$ terminal et avoir une décomposition
  $\varwd = \varwd_1\wdPt \cdots\wdPt \varwd_{n-1}$ où
  $\varwd_i \in \Val(e(\varqP{i},\varqP{i+1}))$.
\end{remark}

On introduit maintenant le lemme essentiel menant au théorème souhaité~:
éliminer un état avec la procédure définie plus tôt ne change pas le langage
calculé.

\begin{lemma}\label{lem.elim.state}
  Soit $\ABA$ un alphabet et $\varAuA$ un automate généralisé contenant
  exactement une transition entre chaque deux états. Soit $\varq$ un état ni
  initial ni terminal. On définit l'automate $\varAuB$ par~:
  \begin{itemize}
  \item $\varQB \defeq \varQA \setminus \{q\}$
  \item $\varQP{0,\varAuB} \defeq \varQP{0,\varAuA}$
  \item $e(\varqP{1},\varqP{2}) \defeq e(\varqP{1},\varq)
    \toSt{(e(\varq,\varq))} e(\varq,\varqP{2}) + e(\varqP{1},\varqP{2})$
  \item $\varAFP{\varAuB} \defeq \varAFP{\varAuA}$
  \end{itemize}
  Alors $\varAuB$ reconnaît le même langage que $\varAuA$.
\end{lemma}

\begin{proof}
  Si on trouve une trace $(\varqZ,\ldots,\varqP{n})$ acceptante pour un mot
  $\varwd$ dans $\varAuA$, alors en supprimant les occurrences de $\varq$ dans
  cette trace on obtient une trace acceptante pour $\varwd$ dans $\varAuB$,
  puisque pour $\varqP{i-1},\varq,\varqP{i+1}$ et deux facteurs
  $\varwd_{i-1},\varwd_i$ apparaissant dans la trace (respectivement la
  décomposition de $\varwd$), on peut les remplacer par
  $\varqP{i},\varqP{i+1}$ et $\varwd_{i-1}\wdPt \varwd_i$, et cela donnera une
  transition valide pour $\varAuB$. Entre deux états $\varqP{i-1},\varqP{i}$ où
  $\varq$ n'apparaît pas, comme on sait que
  $\Val(e_{\varAuA}(\varqP{i-1},\varqP{i}))\subseteq
  \Val(e_{\varAuB}(\varqP{i-1},\varqP{i}))$
  (car on a une union), on en déduit le résultat pour le reste de la trace.

  Réciproquement, si on a une trace valide pour $\varAuB$ alors, entre deux
  états $\varqP{i},\varqP{i+1}$, on a soit un mot $\varwd_i$ qui est dans
  $\Val(e_{\varAuA}(\varqP{i},\varqP{i+1}))$, soit on peut remplacer ces deux
  états par $(\varqP{i},\varq,\varqP{i+1})$ pour faire une trace valide pour
  $\varAuA$.
\end{proof}

\noindent
On peut donc enfin construire une expression rationnelle à partir d'un automate.

\begin{proposition}
  Soit un langage reconnaissable $\varLL$, alors il existe une expression
  rationnelle $e$ telle que $\Val(e) = \varLL$.
\end{proposition}

\begin{proof}
  On sait qu'un langage reconnaissable peut être reconnu par un automate non
  déterministe avec $\Ewd$-transitions. On pose donc $\mathcal A$ un tel
  automate, reconnaissant $\varLL$. On le normalise et lui associe naturellement
  un automate généralisé $\varAuB$ qui contient exactement une transition entre
  deux de ses états. En utilisant le \cref{lem.elim.state}, on sait donc qu'on
  peut éliminer n'importe quel état qui n'est ni initial ni terminal~: on peut
  donc itérer ce lemme pour obtenir un automate généralisé $\varAuA''$ qui
  reconnait le même langage $\varLL$ que $\varAuB$, et qui ne contient qu'un
  état initial et un état terminal (car l'automate a été choisi normalisé).
  Notre automate est donc de la forme~:
  \begin{figure}[h]
    \centering
    \begin{tikzpicture}[node distance = 2cm, on grid, auto]
      \node[state,initial] (q_0) {$\varqZ$};
      \node[state,accepting] (q_1) [right = of q_0] {$\varqP{1}$};
      \path[->,>=latex]
      (q_0) edge node {$e$} (q_1);
    \end{tikzpicture}
  \end{figure}
  
  On en déduit donc que pour tout mot $\varwd \in \toStA$, on a l'équvalence
  suivante~:
  \[\varwd \in \varLL \iff \varwd \in \Val(e)\]
  Donc il existe $e$ tel que $\varLL = \Val(e)$.
\end{proof}

\noindent
On a ainsi les deux sens du théorème de Kleene.

\begin{theorem}[Kleene \cite{Kleene1951RepresentationOE}]
  Soit $\ABA$ un alphabet. Les classes $\Ratio(\ABA)$ et $\Reco(\ABA)$ sont
  égales.
\end{theorem}

\begin{remark}
  Ces deux équivalences sont algorithmiques~: elles permettent de construire
  de manière effective, étant donné un automate $\varAuA$, une expression
  rationnelle $e(\varAuA)$ reconnaissant le même langage, et inversement de
  construire un automate $\varAuA(e)$ depuis une expression rationnelle $e$
  (il faut alors utiliser l'automate union, la concaténation et l'étoile
  d'automates pour obtenir l'automate depuis une expression rationnelle).
\end{remark}

Nous avons vu de nombreuses constructions à propos des automates. Ces
constructions renforcent le fait qu'un langage reconnaissable est reconnu par de
nombreux automates. On peut donc se demander s'il existe un automate meilleur
que les autres, ce qui nous mène à la prochaine section~: la construction d'un
automate minimal.

\`A partir de maintenant, on parlera toujours de langage rationnel pour parler
indistinctement de langage reconnaissable ou rationnel, puisque les deux notions
sont équivalentes.

\section{Vers le théorème de Myhill-Nérode}

Le théorème au c\oe ur de cette section possède deux aspects~: le premier est
la notion d'automate minimal, que nous avons évoqué juste avant. Le deuxième
sera notre premier objectif~: donner un critère permettant de prouver qu'un
langage est (ou non) rationnel. Nous allons aborder cet aspect du théorème en
premier lieu.

\subsection{Conditions de rationnalité d'un langage}

Pour arriver à ce résultat, on va présenter un premier lemme important~: le
lemme de l'étoile. Celui-ci nous indique une condition nécessaire à ce qu'un
langage soit rationnel, mais non suffisante.

\begin{lemma}[de l'étoile]
  Soit un langage rationnel $\varLL$ sur un alphabet $\ABA$. Alors il existe un
  entier $N \in \bN$ tel que pour tout mot $\varwd\in\varLL$ de longueur
  supérieure à $N$, il existe une décomposition $\varwd = \varwx\varwv\varwy$
  telle que
  \[\begin{cases}
  |\varwv| > 0\\
  |\varwx\varwv| \leq N\\
  \forall n \in \bN, \varwx\varwv^n\varwy \in \mathcal L
  \end{cases}\]
\end{lemma}

La preuve de ce résultat se base principalement sur l'idée de trouver une boucle
dans la lecture d'un mot~: si un mot est assez long, alors il finira par
repasser par un état, et on pourra alors itérer cette boucle. En prenant
une trace $(\varqZ,\ldots,\varqP{f})$, passant deux fois par $\varqP{i}$, on
voudra avoir le chemin suivant~:

\begin{figure}[H]
  \centering
  \begin{tikzpicture}[node distance = 2cm, on grid, auto]
      \node[state,initial] (q_0) {$\varqZ$};
      \node[state] (q_1) [right = of q_0] {$\varqP{i}$};
      \node[state] (q_2) [right = of q_1] {$\varqP{i}$};
      \node[state,accepting] (q_3) [right = of q_2] {$\varqP{f}$};
      \path[->,>=latex]
      (q_0) edge node {$\varwx$} (q_1)
      (q_1) edge node {$\varwv$} (q_2)
      (q_2) edge node {$\varwy$} (q_3);
  \end{tikzpicture}
\end{figure}

\begin{proof}
  Puisque $\varLL$ est un langage rationnel, on pose $\varAuA$ un automate
  (déterministe, sans $\Ewd$-transitions) reconnaissant $\varLL$. On pose
  $N = \Card(\varQQ) + 1$, prouvons que la propriété souhaitée est vraie pour ce
  choix de $N$~:

  Soit $\varwd\in\varLL$ de longueur $n \geq N$. Comme $\varwd \in \varLL$, on
  peut trouver une trace acceptante $(\varqZ,\ldots,\varqP{n})$ pour $\varwd$
  dans $\varAuA$. Comme $n \geq N$, on en déduit que $n > |\varQQ|$, donc par le
  principe des tiroirs, on trouve $\varq$ tel qu'il existe au moins deux
  occurrences de $\varq$ dans notre trace. On coupe alors notre trace selon deux
  telles occurrences~: $(\varqZ,\ldots,\varqP{u})$,
  $(\varqP{i},\ldots,\varqP{j})$ et $(\varqP{j},\ldots,\varqP{n})$ où
  $\varqP{i} = \varqP{j} = \varq$. De même, on peut découper
  $\varwd = \varwx\varwv\varwy$ de telle sorte que
  $(\varqZ,\ldots,\varqP{i})$ est une trace pour $\varwx$ de $\varqZ$ à $\varq$,
  $(\varqP{i},\ldots,\varqP{j})$ est une trace pour $\varwv$ de $\varq$ à
  $\varq$ et
  $(\varqP{j},\ldots,\varqP{n})$ est une trace pour $\varwy$ de $\varq$ à
  $\varqP{n}$. Ce raisonnement peut s'appliquer en considérant uniquement les
  $N$ premier élément de la trace totale, donc on peut trouver cette
  décomposition avec $|\varwx\varwv| \leq N$. De plus, puisque $\varq$ a deux
  occurrences distinctes, on sait que $\varwv \neq \Ewd$.
  Enfin, on remarque qu'en concaténant $p$ fois la suite
  $(\varqP{i},\ldots,\varqP{j})$ au milieu, on obtient une trace acceptante pour
  $\varwx\varwv^p\varwy$. Le résultat est donc vérifié pour
  $N = \Card(\varq+1)$.
\end{proof}

\begin{remark}
  Tout d'abord, il est possible que l'automate ne présente aucune boucle. Dans
  ce cas, le langage est fini puisque le nombre de traces possibles l'est. Notre
  lemme nous donne alors comme $N$ un nombre jamais atteint par un mot~: le
  lemme est alors vrai par vacuité.

  Un autre point à relever concerne le fait que, dans notre décomposition, on a
  aussi le fait que $\varwx\varwy \in \varLL$, car supprimer la boucle conserve
  aussi la validité de la trace.

  Enfin, un point plus important~: ce nombre $N$ existe pour un langage $\varLL$
  donné, mais il n'est pas uniforme puisqu'il dépend du choix de l'automate. On
  pourrait considérer le $N$ minimal pour tous les automates possibles, mais
  celui-ci n'a pas de construction explicite \latinexpr{a priori}. On va en fait
  pouvoir montrer que ce $N$ minimal correspond au $N$ associé à un automate
  particulier, et qu'on peut donc choisir un $N$ de façon uniforme en fonction
  de $\varLL$.
\end{remark}

\begin{exercise}
  En utilisant le lemme précédent, montrer que le langage suivant n'est pas
  rationnel~:
  \[\varLL \defeq \compr{0^n1^n}{n \in \bN}\]
  Faire de même pour le langage
  \[\varLL' \defeq \compr{\varwd\in\toSt{\btwo}}{\wul_0 = \wul_1}\]
\end{exercise}

On souhaite maintenant affiner le résultat en cherchant une condition non
seulement nécessaire, mais aussi suffisante, permettant donc de caractériser les
langages rationnels.

Pour y arriver, on a besoin d'une premier définition~: celle de quotient gauche.

\begin{definition}[Quotient gauche]
  Soit $\ABA$ un alphabet et $\varLL$ un langage. Pour tout $\varwd \in \toStA$,
  on définit le quotient gauche (ou résiduel) de $\varLL$ par $\varwd$ comme~:
  \[\luLq\defeq \compr{\varwv \in \toStA}{\varwd\wdPt\varwv \in \varLL}\]
\end{definition}

Le quotient gauche par $\varwd$ peut se voir au niveau des automates~: si on
prend un automate qui reconnaît $\varLL$, la lecture de $\varwd$ depuis l'état
initial mène à un état $\varq$, et l'ensemble $\luLq$ est l'ensemble des mots
acceptés si l'on remplace l'état initial par $\varq$, c'est-à-dire l'ensemble
des mots qui finissent la lecture de $\varwd$ en une lecture acceptante. On
remarque en particulier que $\Ewd \in \luLq$ si et seulement si
$\varwd \in \varLL$.

Cela motive notre construction syntaxique~: les différents $\luLq$ peuvent
naturellement former une structure d'automates, puisque les transitions peuvent
se définir comme étant la transition
$\trA(\luLq,\varla) \defeq \luLP{(\varwd\wdPt\varla)}$. On a aussi défini les
états terminants avec la remarque précédente, et l'état initial est simplement
$\varLL = \luLP{\Ewd}$. Un tel automate, qu'on appelle automate syntaxique,
respecte notre définition d'automate lorsque le nombre d'états est fini, ce qui
nous donne en fait le théorème de Myhill-Nérode.

\begin{theorem}[Myhill-Nérode \cite{Nerode1963-NERLAT-2}]
  Soit $\ABA$ un alphabet et $\varLL$ un langage sur $\ABA$. Le langage $\varLL$
  est rationnel si et seulement si l'ensemble $\compr{\luLq}{\varwd\in\toStA}$
  est fini.
\end{theorem}

\begin{proof}
  Supposons que $\varLL$ est rationnel. On trouve un automate $\varAuA$
  déterministe reconnaissant $\varLL$. On construit alors une surjection de
  l'ensemble des états vers l'ensemble des $\luLq$ pour $\varwd\in \toStA$. Pour
  cela, pour chaque état $\varq \in \varQQ$ on définit l'automate
  $\varAuP{\varq}$ où l'état initial est remplacé par $\varq$. On définit alors
  l'ensemble
  \[f(\varq) \defeq \compr{\varwd\in\toStA}{\varAuP{\varq} \satisf\varwd}\]
  La fonction $f$ est surjective~: pour tout $\varwd$, en prenant
  $\varqP{\varwd} \defeq \trA(\varqZ,\varwd)$, on voit que
  $f(\varqP{\varwd}) = \luLq$. Ainsi, comme $\varQQ$ est fini, on en déduit que
  $\compr{\luLq}{\varwd\in\toStA}$ est lui aussi fini.

  Réciproquement, supposons que $\compr{\luLq}{\varwd\in\toStA}$ est fini. On
  construit alors l'automate syntaxique $\varAuP{\varLL}$~:
  \begin{itemize}
    \begin{minipage}{0.45\textwidth}
    \item $\varQQ \defeq \compr{\luLq}{\varwd\in\toStA}$
    \item $\varqZ \defeq \varLL$
    \end{minipage}
    \begin{minipage}{0.45\textwidth}
    \item $\trA(\luLq,\varla) \defeq \luLP{(\varwd\wdPt\varla)}$
    \item $\varAF \defeq \compr{\varq \in \varQQ}{\Ewd \in \varq}$
    \end{minipage}
  \end{itemize}
  On veut vérifier que cet automate reconnaît bien $\varLL$. Par une induction
  directe sur $\varwd\in\toStA$, on prouve que
  \[\forall \varwd \in \toStA, \trAS(\varqZ,\varwd) = \luLq\]
  donc $\varAuP{\varLL} \satisf\varwd$ si et seulement si $\Ewd \in \luLq$, ce
  qui signifie exactement que $\varwd\in\varLL$ d'après la remarque précédente.
\end{proof}

Il existe donc, pour tout langage rationnel $\varLL$, un automate naturellement
associé.

\begin{figure}[t]
  \centering
  \begin{tikzpicture}[node distance = 3cm, on grid, auto]
    \node[state,initial] (q_0) {$\varqP{s}$};
    \node[state] (q_1) [above = of q_0] {$\varqZ$};
    \node[state] (q_2) [right = of q_0] {$\varqP{1}$};
    \node[state,accepting] (q_3) [right = of q_1] {$\varqP{f}$};
    \node[state] (q_p) [below left = of q_0] {$\varqP{p}$};
    \path[->,>=latex]
    (q_0) edge node {$0$} (q_1)
    (q_0) edge node {$1$} (q_2)
    (q_1) edge node {$0$} (q_3)
    (q_2) edge node {$0$} (q_3)
    (q_3) edge[loop right] node {$0,1$} ()
    (q_p) edge[loop left] node {$0,1$} ()
    (q_1) edge[bend right] node {$1$} (q_p)
    (q_2) edge[bend left] node {$1$} (q_p);
  \end{tikzpicture}

  \vspace{0.2cm}
  \hrule

  \vspace{0.2cm}
  \caption{Automate $\varAuA_2$}
  \label{fig.auto.2}
\end{figure}

\begin{remark}
  L'une des formulations équivalentes de l'axiome du choix est que toute
  surjection $f : X \to Y$ admet une section, c'est-à-dire une fonction
  $g : Y \to X$ telle que $f\circ g = \id_Y$. Pour cette raison, montrer un
  résultat de finitude (et plus généralement de cardinalité) en utilisant une
  surjection est moins fort que d'utiliser une injection, puisqu'on définit
  la subpotence par l'existence d'une injection.

  Dans notre cas, il est donc légitime de se demander s'il n'est pas possible
  de considérer une injection, et la première idée qui vient en tête pour une
  telle injection est alors, étant donné un automate qui reconnait
  $\varLL$~:
  \[\makeFun{f}{\compr{\luLq}{\varwd\in\toStA}}{\varQQ}
            {\clas{\varwd}}{\trAS(\varqZ,\varwd)}\]
  Malheureusement, cette fonction n'est pas bien définie, car deux mots peuvent
  être dans le même résiduel pour $\varLL$ et pour autant être associé à des
  états différents. Un exemple d'une telle situation est l'automate $\varAuA_2$
  donné en \cref{fig.auto.2}.
  
  En effet, les états $\varqZ$ et $\varqP{1}$ sont associés au même résiduel.
  Heureusement, comme dans notre cas $\varQQ$ est fini, il est possible (par
  récurrence sur le cardinal de $\varQQ$) de choisir un état pour chaque
  résiduel, et donc de construire une fonction $f$ comme ci-dessus sur un
  sous-ensemble de $\varQQ$.

  Remarquons aussi que cette fonction $f$ est bien définie exactement lorsque
  l'automate est isomorphe à l'automate syntaxique.
\end{remark}

Cela nous donne un moyen de montrer à la fois qu'un langage est rationnel, et
qu'un langage n'est pas rationnel~: en exibant une infinité de résiduels, on en
déduit directement qu'un langage n'est pas rationnel, et ce critère est plus
précis que le lemme de l'étoile puisqu'il est exact.

\begin{exercise}
  On considère l'alphabet $\{a,b,c\}$, montrer que le langage
  \[\varLL \defeq \compr{a^nb^mc^p}{(n,m,p)\in (\bN^*)^3}\]
  vérifie le lemme de l'étoile, mais n'est pas rationnel.
\end{exercise}

\subsection{Automate minimal}

On passe maintenant à l'autre aspect du théorème de Myhill-Nérode~: l'automate
syntaxique qu'on a construit est en fait l'automate minimal, au sens où tout
automate reconnaissant le même langage contiendra plus d'états.

Pour en donner une idée plus claire, on va introduire les notions de morphismes
d'automates et d'automate quotient. Pour simplifier, on ne parlera que
d'automates finis déterministes complets et accessibles.

\begin{definition}[Morphisme d'automate]
  Soit $\ABA$ et $\varAuA$, $\varAuB$ deux automates finis déterminstes sur
  $\ABA$. On appelle morphisme d'automate $\makeFunName{f}{\varAuA}{\varAuB}$
  une fonction $\makeFunName{f}{\varQA}{\varQB}$ telle que~:
  \begin{itemize}
  \item $f(\varqP{0,\varAuA}) = \varqP{0,\varAuB}$
  \item pour tous $\varq \in \varQA, \varla \in \ABA$,
    $f(\trAA(\varq,\varla)) = \trAB(f(\varq),\varla)$
  \item pour tout $\varq \in \varAAF, f(q) \in \varABF$
  \end{itemize}
\end{definition}

Un morphisme d'automate est donc une façon de transformer l'automate $\varAuA$
en l'automate $\varAuB$. Remarquons que généralement, il est tout à fait
possible de n'avoir des morphismes que dans un sens, par exemple entre
l'automate $\varAuA_3$ et l'automate $\varAuA_4$, donnés en \cref{fig.auto.3}.
On peut trouver un morphisme de l'automate de gauche vers l'automate de droite,
en prenant $f(\varqP{1}) = f(\varqP{2}) = \varqP{1}'$. Cependant, dans l'autre
sens, on ne peut pas conserver la transition $0,1$ entre $\varqZ'$ et
$\varqP{1}'$ en associant une image à $\varqP{1}'$.

\begin{figure}[t]
  \centering
  \begin{tikzpicture}[node distance = 2cm, on grid, auto]
      \node[state,initial] (q_0) {$\varqZ$};
      \node[state,accepting] (q_1) [above right = of q_0] {$\varqP{1}$};
      \node[state, accepting] (q_2) [below right = of q_0] {$\varqP{2}$};
      \path[->,>=latex]
      (q_0) edge node {$0$} (q_1)
      (q_0) edge node {$1$} (q_2)
      (q_1) edge [loop right] node {$0,1$} ()
      (q_2) edge [loop right] node {$0,1$} ();
      \node (use) [right = of q_0] {};
      \node (use2) [right = of use] {};
      \node[state, initial] (q_3) [right = of use2] {$\varqZ'$};
      \node[state, accepting] (q_4) [right = of q_3] {$\varqP{1}'$};
      \path[->,>=latex]
      (q_3) edge node {$0,1$} (q_4)
      (q_4) edge [loop right] node {$0,1$} ();
  \end{tikzpicture}

  \vspace{0.2cm}
  \hrule

  \vspace{0.2cm}
  \caption{Automates $\varAuA_3$ et $\varAuA_4$}
  \label{fig.auto.3}
\end{figure}

Un premier résultat nous dit qu'au niveau des langages reconnus, un morphisme
donne une inclusion.

\begin{proposition}
  Soient $\varAuA, \varAuB$ deux automates sur un alphabet $\ABA$, et
  $\makeFunName{f}{\varAuA}{\varAuB}$ un morphisme d'automates. Alors
  \[\varLL_{\varAuA} \subseteq \varLL_{\varAuB}\]
\end{proposition}

\begin{proof}
  On va noter $\varq,\varqZ,\ldots$ les états de $\varAuA$ et
  $\varq',\varqZ',\ldots$ les états de $\varAuB$.
  
  On voit tout d'abord que pour tout mot $\varwd\in\toStA$, on a
  \[\trAS(f(\varqZ),\varwd) = f(\trAS(\varqZ,\varwd))\]
  et $f(\varqZ) = \varqZ'$. De plus, comme $f$ est un morphisme,
  $f(\varAAF)\subseteq \varABF$, donc
  \[\trAS(\varqZ,\varwd)\in\varAAF\implies f(\trAS(\varqZ,\varwd))\in\varABF
  \qedhere\]
\end{proof}

Un autre résultat, peut-être moins intuitif, est que l'inclusion devient une
égalité dès lors que $f(\varAAF) = \varABF$.

\begin{proposition}
  Si un morphisme d'automates $\makeFunName{f}{\varAuA}{\varAuB}$ est tel que
  $f(\varAAF) = \varABF$, alors $\varLL_{\varAuA} = \varLL_{\varAuB}$.
\end{proposition}

\begin{proof}
  On réutilise les notations de la preuve précédente. Par contraposée, supposons
  que $\varwd \notin \varLL_{\varAuA}$. Cela signifie donc que
  $\trAS(\varqZ,\varwd) \notin \varAAF$, donc
  $f(\trAS(\varqZ,\varwd)) \notin \varABF$ puisque
  $\varABF = f(\varAAF)$, donc $\trAS(\varqZ',\varwd)\notin \varABF$ par
  l'argument précédent.

  Par contraposée, $\varAuB\satisf \varwd \implies \varAuA \satisf \varwd$, d'où
  le résultat.
\end{proof}

On dit qu'un tel morphisme est surjectif (attention, cela ne signifie pas que le
morphisme est surjectif en tant que fonction $\varQA \to \varQB$).
On peut alors prouver que l'automate syntaxique est minimal en un sens précis.

\begin{proposition}
  Soit $\varAuA$ un automate fini déterministe complet et accessible sur un
  alphabet $\ABA$. Il existe un morphisme surjectif
  $\varAuA \to \varAuA_{\varLL_{\varAuA}}$ qui est aussi surjectif sur les états.
\end{proposition}

\begin{proof}
  Comme on sait que $\varAuA$ est accessible, on peut trouver pour tout état
  $\varq \in \varQQ$ un mot $\varwd(\varq)$ tel que
  $\trAS(\varqZ,u(\varq)) = \varq$. On définit alors l'image de $\varq$ comme
  $\lQuot{(\varwd(\varq))}{\varLL_{\varAuA}}$. On montre d'abord que
  ce choix d'image ne dépend pas du mot $\varwd(\varq)$ choisi. En effet, si
  $\trAS(\varqZ,\varwd) = \trAS(\varqZ,\varwv)$ alors pour tout
  $\varww \in \toStA$~:
  \[\varAuA \satisf \varwd\wdPt\varww \iff \varAuA \satisf\varwv\wdPt\varww\]
  donc $\lQuot{\varwd}{\varLL_{\varAuA}}= \lQuot{\varwv}{\varLL_{\varAuA}}$.

  Pour un état terminant dans l'automate syntaxique donné, on trouve un mot
  $\varwd$ tel que cet état vaut $\lQuot{\varwd}{\varLL_{\varAuA}}$. On peut
  alors considérer $\trAS(\varqZ,\varwd)$ pour obtenir un état qui est associé à
  cet état, et puisque $\Ewd \in \lQuot{\varwd}{\varLL_{\varAuA}}$, on en déduit
  que $\trAS(\varqZ,\varwd)\in \varAAF$, donc notre morphisme est surjectif. Il
  est même surjectif en tant que fonction entre états, puisque l'argument
  précédent donne pour tout $\lQuot{\varwd}{\varLL_{\varAuA}}$ un antécédent par
  le morphisme~: $\trAS(\varqZ,\varwd)$.
\end{proof}

\begin{remark}
  En soi, le point important est la surjectivité sur les états. Cependant, comme
  nous allons le voir, l'existence d'un morphisme nous permet de montrer que
  l'automate minimal peut être obtenu à partir des autres automates, en
  regroupant les états.
\end{remark}

Nous cherchons maintenant à obtenir une construction explicite de l'automate
syntaxique à partir d'un automate donné. Pour cela, nous introduisons d'abord la
notion d'automate quotient. Ces quotients s'appliquent à ce qu'on va appeler des
congruences d'automates, qui sont des relations d'équivalence compatibles avec
la structure d'automate.

\begin{definition}[Congruence d'automate]
  Soit $\ABA$ un alphabet et $\varAuA$ un automate fini déterministe sur $\ABA$.
  Une relation $\mathord{\sim} \subseteq \varQQ \times \varQQ$ est appelée une
  congruence d'automate si toutes les conditions suivantes sont vérifiées~:
  \begin{itemize}
  \item $\sim$ est une relation d'équivalence.
  \item $\sim$ est compatible avec $\trA$~:
    \[\forall \varq,\varq' \in \varQQ, \varq \sim \varq' \implies
    \forall \varla \in \ABA, \trA(\varq,\varla) \sim \trA(\varq',\varla)\]
  \item $\sim$ sature $\varAF$~:
    \[\forall \varq,\varq' \in \varQQ, \varq \sim \varq' \implies
    (\varq \in \varAF \iff \varq' \in \varAF)\]
  \end{itemize}

  On peut alors définir l'automate quotient $\quot{\varAuA}{\sim}$ par~:
  \begin{itemize}
  \item son ensemble d'états est $\quot{\varQQ}{\sim}$.
  \item son état initial est la classe de $\varqZ$.
  \item pour tout $\varq \in \varQQ$ et $\varla \in \ABA$,
    $\trA(\clas{\varq},\varla)=\clas{\trA(\varq,\varla)}$
  \item l'ensemble des états terminaux est
    $\compr{\clas{\varqP{f}}}{\varqP{f} \in F}$.
  \end{itemize}
\end{definition}

\begin{remark}
  La compatibilité est essentielle pour que cette définition ait du sens,
  puisque l'on veut que le choix de $\varq$ dans une classe ne change pas le
  résultat de $\trA(\varq,\varla)$~: c'est le cas ici puisque si
  $\varq\sim \varq'$ alors $\trA(\varq,\varla) \sim \trA(\varq',\varla)$. Pour
  les états terminaux, dire qu'un ensemble d'états est terminal semble changer
  le sens de l'automate lorsque l'ensemble n'est pas constitué que d'états
  terminaux, ce qui motive notre propriété de saturation (remarquons que la
  condition de saturation aurait pu se limiter à une implication, puisque la
  relation est symétrique dans tous les cas).
\end{remark}

On prouve alors que le quotient, comme attendu puisqu'on regroupe des états de
même comportement, ne change pas le langage reconnu.

\begin{proposition}
  Soit un automate $\varAuA$ et une congruence d'automate sur $\varAuA$.
  Alors
  \[\varLL_{\varAuA} = \varLL_{\quot{\varAuA}{\sim}}\]
\end{proposition}

\begin{proof}
  Par une récurrence sur $\varwd$ qu'on ne détaille pas, on vérifie que
  \[\forall \varwd \in \toStA,
  \trASP{\quot{\varAuA}{\sim}}(\clas{\varqZ},\varwd)=
  \clas{\trAAS(\varqZ,\varwd)}\]
  Si $\trAAS(\varqZ,\varwd) \in \varAAF$ alors
  $\clas{\trAAS(\varqZ,\varwd)} \in \varAFP{\quot{\varAuA}{\sim}}$.
  Réciproquement, si
  $\clas{\trAAS(\varqZ,\varwd)} \in \varAFP{\quot{\varAuA}{\sim}}$, alors on
  sait qu'il existe $\varq$ dans cette classe tel que $\varq \in \varAAF$, et
  comme $\sim$ sature $\varAAF$, on en déduit que
  $\trAA(\varqZ,\varwd) \in \varAAF$.

  On a donc l'égalité des deux langages.
\end{proof}

On peut en fait relier directement les morphismes et les congruences~: comme
en algèbre, un morphisme donne lieu à une congruence par laquelle on peut
quotienter notre objet pour obtenir un objet plus petit, et plus \og exact\fg
vis à vis de ce morphisme.

\begin{proposition}
  Soient $\varAuA, \varAuB$ deux automates sur le même alphabet $\ABA$ et
  $\makeFunName{f}{\varAuA}{\varAuB}$ un morphisme surjectif. Alors $f$ définit
  sur $\varAuA$ une congruence d'automate par~:
  \[\varq\sim_f \varq' \defeq f(\varq) = f(\varq')\]
\end{proposition}

\begin{proof}
  On veut montrer que $\sim_f$ définit une congruence. Il est clair que c'est
  une relation d'équivalence, on prouve donc les deux autres points~:
  \begin{itemize}
  \item si $f(\varq) = f(\varq')$ et $\varla \in \ABA$, alors
    \begin{align*}
      f(\trA(\varq,\varla)) &= \trA(f(\varq),\varla)\\
      &= \trA(f(\varq'),\varla)\\
      f(\trA(\varq,\varla)) &= f(\trA(\varq',\varla))
    \end{align*}
    donc $\sim_f$ est compatible avec $\trA$.
  \item si $f(\varq) = f(\varq')$ et $\varq \in \varAF$, alors
    $f(\varq) \in f(\varAF)$, donc $f(\varq') \in f(\varAF)$, d'où par
    surjectivité de $f$ que $\varq' \in \varAF$.
  \end{itemize}
  Ainsi $\sim_f$ est une congruence.
\end{proof}

Vu la forme de $\sim_f$, il est clair que $\quot{\varAuA}{\sim_f}$ est isomorphe
à une partie de $\varAuB$ (en prenant pour notion d'isomorphisme l'existence
d'un morphisme réciproque). On en déduit donc qu'on peut construire
$\varAuP{\varLL_{\varAuA}}$ directement depuis $\varAuA$ par un quotient~: il se
trouve que ce quotient peut être décrit de façon très concrète.

\begin{definition}[Congruence de Nérode \cite{Nerode1963-NERLAT-2}]
  Soit $\varAuA$ un automate sur un alphabet $\ABA$. On définit la congruence de
  Nérode sur $\varAuA$ par~:
  \[\varq \equiv \varq' \defeq \forall \varwd \in \toStA,
  \trAS(\varq,\varwd) \in \varAF \iff \trAS(\varq',\varwd) \in \varAF\]
\end{definition}

On voit déjà que si on prend deux mots $\varwd$ et $\varwv$ tels que
$\trAS(\varqZ,\varwd)=\varq$ et $\trAS(\varqZ,\varwv) = \varq'$, alors
$\varq \equiv \varq'$ revient à dire que
$\lQuot{\varwd}{\varLL_{\varAuA}}= \lQuot{\varwv}{\varLL_{\varAuA}}$, ce qui
confirme que le quotient de $\varAuA$ par $\equiv$ nous donne l'automate
syntaxique. Pour l'instant, notre construction de l'automate ne peut pas être
automatisée~: on a besoin de tester tous les mots pour déterminer si
$\varq\sim \varq'$. Pour combler cette lacune, on définit les approximations de
la congruence de Nérode~:

\begin{definition}[Approximation de la congruence de Nérode]
  On définit l'approximation $n$ de la congruence de Nérode, notée $\equiv_n$,
  par
  \[\varq\equiv_n \varq' \defeq \forall \varwd \in \ABA^{<n},
  \trAS(\varq,\varwd)\in \varAF \iff \trAS(\varq',\varwd) \in \varAF\]
\end{definition}

L'approximation consiste donc en la considération des premiers mots seulement.
On introduit maintenant un lemme permettant de calculer effectivement cette
congruence pour tout $n$.

\begin{lemma}\label{lem.nerode.succ}
  Soit un automate $\varAuA$ sur un alphabet $\ABA$, deux états $\varq,\varq'$
  de $\varAuA$ et $n \in \bN$. On a la relation suivante~:
  \[\varq \equiv_{n+1} \varq' \iff \varq\equiv_n \varq' \land
  \forall \varla \in \ABA,
  \trA(\varq,\varla)\equiv_n \trA(\varq',\varla)\]
\end{lemma}

\begin{proof}
  Si $\varq\equiv_{n+1}\varq'$, alors pour tout $\varla\in \ABA$ et
  $\varwd \in \ABA^{<n}$, on a l'équivalence
  $\trAS(\varq,\varwd\wdPt \varla) \in \varAF \iff
  \trAS(\varq',\varwd\wdPt \varla) \in \varAF$
  par hypothèse. Mais alors, cela se réécrit en
  \[\trAS(\trA(\varq,\varla),\varwd)\in\varAF \iff
  \trAS(\trA(\varq',\varla),\varwd)\in\varAF\]
  On en déduit que $\trA(\varq,\varla)\equiv_n \trA(\varq',\varla)$.
  Réciproquement, pour $\varwd \in \ABA^{< n + 1}$, on peut éliminer le cas où
  $\varwd = \Ewd$ qui est direct et considérer que $\varwd = \varwv\wdPt\varla$
  où $\wvl < n$, ce qui nous permet d'utiliser l'hypothèse pour déduire que
  \[\trAS(\varq,\varwv\wdPt \varla)\in \varAF \iff
  \trAS(\varq',\varwv\wdPt \varla)\]
  c'est-à-dire $\varq\equiv_{n+1} \varq'$.
\end{proof}

Le deuxième lemme nous dit que l'algorithme effectué en construisant cette suite
de congruence, lorsqu'il atteint un point stationnaire, donne en résultat la
congruence de Nérode.

\begin{lemma}
  Soit un automate $\varAuA$ sur un alphabet $\ABA$. Si pour $n \in \bN$, les
  relations $\equiv_n$ et $\equiv_{n+1}$ coïncident, alors pour tout $k \in \bN$,
  $\equiv_{n+k}$ et $\equiv$ coïncident.
\end{lemma}

\begin{proof}
  Supposons que pour tous $\varq,\varq'$, on ait
  \[\varq\equiv_n \varq' \iff \varq\equiv_{n+1} \varq'\]
  alors par induction sur $k$, on prouve que $\equiv_{n+k}$ coïncide avec
  $\equiv_n$~:
  \begin{itemize}
  \item c'est évident pour $k = 0$.
  \item supposons que $\equiv_{n+k}$ et $\equiv_n$ coïncident. Alors, par le
    \cref{lem.nerode.succ}, on voit que $\equiv_{n+k+1}$ et $\equiv_{n+1}$
    coïncident, mais par hypothèse $\equiv_{n+1}$ et $\equiv_n$ coïncident,
    d'où le fait que $\equiv_{n+1}$ et $\equiv_n$ coïncident.
  \end{itemize}

  On veut maintenant prouver que $\equiv_n$ et $\equiv$ coïncident. Soit deux
  états $\varq,\varq'$ tels que $\varq\equiv_n \varq'$ et un mot
  $\varwd \in \toStA$. Si $\wul < n$ alors par définition
  $\trAS(\varq,\varwd) \in \varAF \iff \trAS(\varq',\varwd) \in \varAF$.
  Sinon, alors $\varq\equiv_{\wul} \varq'$ puisque $\equiv_n$ coïncide avec
  $\equiv_k$ pour tout $k > n$, d'où le résultat. La réciproque est immédiate.
\end{proof}

On a donc un algorithme permettant de calculer l'automate minimal d'un automate
complet et accessible~: on construit la congruence de Nérode par approximations
successives grâce au \cref{lem.nerode.succ} jusqu'à obtenir une partition stable
des états, et le quotient est alors directement l'automate minimal.

\subsection{Les langages reconnus par monoïde}

Pour clore ce chapitre, revenons aux considérations algébriques énoncées au
début. Dans celles-ci, on a vu qu'il était possible d'associer un langage à un
monoïde en utilisant la propriété universelle du monoïde libre. Nous allons
maintenant prouver qu'en réalité, les langages associés à des monoïdes finis
sont exactement les langages rationnels.

On commence par prouver un premier sens~: les langages reconnus par monoïdes
sont aussi reconnus par automates.

\begin{proposition}
  Soit $\ABA$ un alphabet, $\varMMo$ un monoïde et $\varLL$ un langage reconnu
  par une fonction $h : \ABA \to \varMMo$ et $S\subseteq \varMMo$. Alors
  $\varLL \in \Reco(\ABA)$.
\end{proposition}

\begin{proof}
  On construit l'automate $\varAuA$ suivant~:
  \begin{itemize}
    \begin{minipage}{0.45\textwidth}
    \item $\varQQ \defeq \varMMo$
    \item $\varqZ \defeq \neuMM$
    \end{minipage}
    \begin{minipage}{0.45\textwidth}
    \item $\trA(\varq,\varla) \defeq \varq\opM f(\varla)$
    \item $\varAF \defeq S$
    \end{minipage}
  \end{itemize}

  Une récurrence immédiate sur $\varwd$ permet de prouver que
  \[\forall\varwd\in\toStA,\trAS(\varqZ,\varwd)=\prod_{i = 0}^{\wul-1}\varwd_i\]
  Cette valeur est exactement celle de $\relev{h}(\varwd)$, et être reconnu par
  $\varAuA$ correspond exactement à appartenir à $S$, d'où l'égalité de langage.
\end{proof}

Il nous reste donc à construire un monoïde à partir d'un automate. Pour cela,
on va utiliser la construction du monoïde des transitions.

\begin{definition}[Monoïde des transitions]
  Soit $\ABA$ un alphabet et $\varAuA$ un automate fini déterministe
  complet reconnaissant un langage $\varLL$. On définit le monoïde des
  transitions de $\varAuA$, qu'on notera $\varMMo_{\varAuA}$, par~:
  \begin{itemize}
  \item les éléments du monoïde sont les relations binaires
    $\mathord{\varRR}\subseteq \varQQ\times \varQQ$.
  \item la multiplication est la composition de relations~:
    \[\mathord{\varRR}\neuMM \mathord{\varRR'} \defeq
    \compr{(x,z)}{\exists y \in \varQQ, x\varRR y\land y\varRR z}\]
  \item l'élément neutre est la relation
    \[\Delta_{\varQQ} \defeq \compr{(\varq,\varq)}{\varq \in \varQQ}\]
  \end{itemize}
\end{definition}

\begin{exercise}
  Prouver que $\varMMo_{\varAuA}$ est bien un monoïde.
\end{exercise}

Les éléments du monoïde sont les relations binaires. Cependant, les éléments
d'importances sont plus précisément les relations telles que pour tout couple
$(\varq,\varq')$, on a un chemin de $\varq$ vers $\varq'$ dans $\varAuA$. On
peut donc simplement prendre la relation
$(\varq,\varq') \defeq \trA(\varq,\varla) = \varq'$ paramétrée par $\varla$ pour
obtenir à travers $\trAS$ les relations contenant les chemins entre états
paramétrés par un mot $\varwd$ donné. Un mot est acceptant quand l'état
$\varqZ$ est relié à un état terminal, donc le choix de $S$ est celui des
relations incluant un couple $(\varqZ,\varqP{f})$ où $\varqP{f} \in \varAF$. On
en vient donc au résultat à proprement parler.

\begin{proposition}
  Soit un automate fini déterministe complet $\varAuA$ reconnaissant un
  langage $\varLL$ sur un alphabet $\ABA$, alors la fonction
  \[\makeFun{h}{\ABA}{\varMMo_{\varAuA}}{\varla}
            {\compr{(\varq,\varq')}{\trA(\varq,\varla) = \varq'}}\]
  et la partie
  \[S \defeq \compr{\mathord{\varRR} \in \varMMo_{\varAuA}}{
    \mathord{\varRR} \cap (\{\varqZ\}\times \varAF)\neq\vide}\]
  reconnaissent le langage $\varLL$.
\end{proposition}

\begin{proof}
  Par induction sur la longueur de $\varwd$, on montre que~:
  \[\forall \varwd\in \toStA, \relev{h}(\varwd) =
  \compr{(\varq,\varq')}{\trAS(\varq,\varwd) = \varq'}\]
  \begin{itemize}
  \item on sait que $\relev{h}(\Ewd)$ est l'élément neutre du monoïde,
    et que cet élément est $\{(\varq,\varq)\}$, or c'est aussi l'ensemble décrit
    par $\trAS(\varq,\Ewd) = \varq'$.
  \item supposons que l'égalité est vraie pour un mot $\varwd$, et soit
    $\varla \in \ABA$. Alors
    \begin{align*}
      \relev{h}(\varwd\wdPt\varla)
      &= \relev{h}(\varwd) \opM
      \compr{(\varq,\varq')}{\trA(\varq,\varla) = \varq'}\\
      &= \compr{(\varq,\varq')}{\trAS(\varq,\varwd) = \varq'}\opM
      \compr{(\varq,\varq')}{\trA(\varq,\varla) = \varq'}\\
      &= \compr{(\varq,\varq'')}{\exists \varq',
      \trAS(\varq,\varwd) = \varq' \land \trA(\varq',\varla) = \varq''}\\
      &= \compr{(\varq,\varq'')}{\trA(\trAS(\varq,\varwd),\varla) = \varq''}\\
      &= \compr{(\varq,\varq')}{\trAS(\varq,\varwd\wdPt \varla) = \varq'}
    \end{align*}
  \end{itemize}
  D'où le résultat par induction. On sait donc que $\relev{h}(\varwd) \in S$ est
  équivalent à ce qu'il existe un couple
  $(\varqZ,\varqP{f}) \in \compr{(\varq,\varq')}{\trAS(\varq,\varwd) = \varq'}$
  où $\varqP{f} \in \varAF$, ce qui revient à dire que
  $\trAS(\varqZ,\varwd) \in \varAF$, donc à ce que $\varAuA$ accepte $\varwd$.
  Ainsi $\varMMo_{\varAuA}$ reconnaît le même langage que $\varAuA$.
\end{proof}
